% Recovered from the previously completed PDF; no new mathematical content.
% Prose is editable; formula boxes retain original symbol positions.
% Build from repository root: xelatex -no-shell-escape -halt-on-error -output-directory=... print/fudan/NAME.tex
\documentclass[11pt,a4paper]{article}
\usepackage[margin=23mm,headheight=14pt,headsep=8.5mm,footskip=10mm]{geometry}
\usepackage{amsmath,amssymb}
\DeclareMathAlphabet{\mathsf}{OT1}{cmss}{m}{n}
\usepackage[no-math]{fontspec}
\usepackage{xeCJK}
\setmainfont{Latin Modern Roman}
\setCJKmainfont{FandolSong-Regular.otf}[Path=assets/fonts/fandol/,BoldFont=FandolSong-Bold.otf,ItalicFont=FandolSong-Regular.otf,Script=CustomDefault]
\xeCJKsetup{PunctStyle=quanjiao,CheckSingle=true}
\usepackage{setspace}\setstretch{1.12}
\setlength{\parindent}{2em}\setlength{\parskip}{3pt}
\emergencystretch=6em
\hfuzz=9pt
\usepackage{titlesec}
\titleformat{\section}{\centering\normalfont\Large\bfseries}{}{0pt}{}
\titlespacing*{\section}{0pt}{1.8\baselineskip}{1.0\baselineskip}
\usepackage{fancyhdr}\pagestyle{fancy}\fancyhf{}
\fancyhead[L]{\small 数学手记·高等代数}\fancyhead[R]{\small 22S02–22S08}
\fancyfoot[C]{\thepage}\renewcommand{\headrulewidth}{0.4pt}
\usepackage[hidelinks,unicode]{hyperref}
\hypersetup{pdftitle={22S02–22S08：题目、解答与小结},pdfauthor={Vincent Arqain}}
\newcommand{\RB}[4]{\raisebox{0pt}[#2bp][#3bp]{\hbox to #1bp{#4\hss}}}
\newcommand{\RG}[5]{\hbox to 0pt{\kern#1bp\raisebox{#2bp}[0pt][0pt]{\fontsize{#3bp}{#3bp}\csname font#4\endcsname\char#5}\hss}}
\newcommand{\RT}[4]{\hbox to 0pt{\kern#1bp\raisebox{#2bp}[0pt][0pt]{\fontsize{#3bp}{#3bp}\selectfont#4}\hss}}
\newcommand{\RS}[4]{\hbox to 0pt{\kern#1bp\raisebox{#2bp}[0pt][0pt]{\fontsize{#3bp}{#3bp}\selectfont\ensuremath{#4}}\hss}}
\newcommand{\RR}[4]{\hbox to 0pt{\kern#1bp\raisebox{#2bp}[0pt][0pt]{\rule{#3bp}{#4bp}}\hss}}
\newcommand{\fontmi}{\usefont{OML}{cmm}{m}{it}}
\newcommand{\fontsy}{\usefont{OMS}{cmsy}{m}{n}}
\newcommand{\fontex}{\usefont{OMX}{cmex}{m}{n}}
\newcommand{\fontrm}{\usefont{OT1}{cmr}{m}{n}}
\newcommand{\fontbf}{\usefont{OT1}{cmr}{bx}{n}}
\newcommand{\fontam}{\usefont{U}{msa}{m}{n}}
\newcommand{\fontbb}{\usefont{U}{msb}{m}{n}}
\widowpenalty=10000\clubpenalty=10000\raggedbottom
\begin{document}
\begin{center}{\LARGE\bfseries 22S02–22S08\par}\vspace{8pt}复旦大学高等代数习题集·题目与解答\end{center}
% Original page 1.
\section*{2022S02}
求下列\allowbreak{}\hskip.1em \RB{24.941}{7.763}{1.408}{%
\RG{0.000}{-0.003}{10.9091}{mi}{110}\RG{8.972}{-0.003}{10.9091}{rm}{43}\RG{19.870}{-0.003}{10.9091}{rm}{49}%
}\hskip.1em\allowbreak{} 阶方阵\allowbreak{}\hskip.1em \RB{8.366}{8.308}{0.503}{%
\RG{0.000}{-0.003}{10.9091}{mi}{65}%
}\hskip.1em\allowbreak{} 的特征值和特征向量：

\[
\RB{171.070}{26.410}{59.661}{%
\RG{22.727}{25.484}{10.9092}{ex}{48}\RG{157.111}{25.484}{10.9092}{ex}{49}\RG{32.272}{15.171}{10.9091}{mi}{110}\RG{53.935}{15.171}{10.9091}{sy}{0}\RG{62.420}{15.171}{10.9091}{mi}{110}\RG{22.727}{6.284}{10.9092}{ex}{66}\RG{157.111}{6.284}{10.9092}{ex}{67}\RG{32.819}{-0.004}{10.9091}{rm}{49}\RG{48.783}{-0.004}{10.9091}{mi}{110}\RG{57.755}{-0.004}{10.9091}{sy}{0}\RG{68.665}{-0.004}{10.9091}{rm}{50}\RG{84.079}{-0.004}{10.9091}{rm}{49}\RG{91.961}{-0.004}{10.9091}{sy}{0}\RG{102.870}{-0.004}{10.9091}{mi}{110}\RG{22.727}{-0.262}{10.9092}{ex}{66}\RG{157.111}{-0.262}{10.9092}{ex}{67}\RG{22.727}{-6.807}{10.9092}{ex}{66}\RG{157.111}{-6.807}{10.9092}{ex}{67}\RT{119.987}{-12.687}{10.9091}{.}\RG{22.727}{-13.353}{10.9092}{ex}{66}\RG{157.111}{-13.353}{10.9092}{ex}{67}\RT{124.232}{-15.675}{10.9091}{.}\RT{128.477}{-18.664}{10.9091}{.}\RG{0.000}{-19.353}{10.9091}{mi}{65}\RG{11.212}{-19.353}{10.9091}{rm}{61}\RG{168.475}{-19.353}{10.9091}{mi}{58}\RG{58.724}{-19.660}{10.9091}{rm}{50}\RG{84.083}{-19.660}{10.9091}{mi}{110}\RG{93.055}{-19.660}{10.9091}{sy}{0}\RG{103.964}{-19.660}{10.9091}{rm}{52}\RG{22.727}{-19.898}{10.9092}{ex}{66}\RG{157.111}{-19.898}{10.9092}{ex}{67}\RG{22.727}{-26.444}{10.9092}{ex}{66}\RG{157.111}{-26.444}{10.9092}{ex}{67}\RT{90.989}{-32.343}{10.9091}{.}\RT{119.987}{-32.343}{10.9091}{.}\RG{22.727}{-32.989}{10.9092}{ex}{66}\RG{157.111}{-32.989}{10.9092}{ex}{67}\RT{95.234}{-35.332}{10.9091}{.}\RT{124.232}{-35.332}{10.9091}{.}\RT{99.479}{-38.321}{10.9091}{.}\RT{128.477}{-38.321}{10.9091}{.}\RG{142.625}{-39.317}{10.9091}{sy}{0}\RG{151.110}{-39.317}{10.9091}{rm}{49}\RG{22.727}{-39.971}{10.9092}{ex}{64}\RG{157.111}{-39.971}{10.9092}{ex}{65}\RG{122.474}{-54.492}{10.9091}{mi}{110}\RG{142.078}{-54.492}{10.9091}{sy}{0}\RG{150.563}{-54.492}{10.9091}{mi}{110}%
}
\]
解：以\allowbreak{}\hskip.1em \RB{46.684}{7.764}{2.607}{%
\RG{0.000}{-0.002}{10.9091}{rm}{48}\RG{5.454}{-0.002}{10.9091}{mi}{59}\RG{10.303}{-0.002}{10.9091}{rm}{49}\RG{15.758}{-0.002}{10.9091}{mi}{59}\RG{20.591}{-0.002}{10.9091}{mi}{58}\RG{25.435}{-0.002}{10.9091}{mi}{58}\RG{30.267}{-0.002}{10.9091}{mi}{58}\RG{35.111}{-0.002}{10.9091}{mi}{59}\RG{39.955}{-0.002}{10.9091}{mi}{110}%
}\hskip.1em\allowbreak{} 为行列指标，则

\[
\RB{228.424}{7.770}{3.759}{%
\RG{0.000}{0.004}{10.9091}{mi}{97}\RG{18.944}{0.004}{10.9091}{rm}{61}\RG{30.460}{0.004}{10.9091}{mi}{110}\RG{39.432}{0.004}{10.9091}{sy}{0}\RG{50.341}{0.004}{10.9091}{rm}{50}\RG{55.795}{0.004}{10.9091}{mi}{106}\RG{60.290}{0.004}{10.9091}{mi}{59}\RG{86.952}{0.004}{10.9091}{mi}{97}\RG{116.739}{0.004}{10.9091}{rm}{61}\RG{128.254}{0.004}{10.9091}{mi}{106}\RG{132.749}{0.004}{10.9091}{mi}{59}\RG{159.411}{0.004}{10.9091}{mi}{97}\RG{189.198}{0.004}{10.9091}{rm}{61}\RG{200.713}{0.004}{10.9091}{mi}{106}\RG{208.255}{0.004}{10.9091}{sy}{0}\RG{219.164}{0.004}{10.9091}{mi}{110}\RG{225.709}{0.004}{10.9091}{mi}{59}\RG{5.766}{-1.633}{7.9701}{mi}{106}\RG{9.177}{-1.633}{7.9701}{mi}{59}\RG{11.528}{-1.633}{7.9701}{mi}{106}\RG{92.740}{-1.633}{7.9701}{mi}{106}\RG{96.151}{-1.633}{7.9701}{mi}{59}\RG{98.502}{-1.633}{7.9701}{mi}{106}\RG{102.390}{-1.633}{7.9701}{sy}{0}\RG{108.977}{-1.633}{7.9701}{rm}{49}\RG{165.198}{-1.633}{7.9701}{mi}{106}\RG{168.610}{-1.633}{7.9701}{mi}{59}\RG{170.961}{-1.633}{7.9701}{mi}{106}\RG{174.848}{-1.633}{7.9701}{rm}{43}\RG{181.432}{-1.633}{7.9701}{rm}{49}%
}
\]
其中只保留指标在\allowbreak{}\hskip.1em \RB{36.381}{7.766}{2.605}{%
\RG{0.000}{0.000}{10.9091}{rm}{48}\RG{5.455}{0.000}{10.9091}{mi}{59}\RG{10.288}{0.000}{10.9091}{mi}{58}\RG{15.132}{0.000}{10.9091}{mi}{58}\RG{19.964}{0.000}{10.9091}{mi}{58}\RG{24.808}{0.000}{10.9091}{mi}{59}\RG{29.652}{0.000}{10.9091}{mi}{110}%
}\hskip.1em\allowbreak{} 内的元素。

先在多项式空间\allowbreak{}\hskip.1em \RB{32.139}{8.678}{3.495}{%
\RG{0.000}{-0.004}{10.9091}{bb}{67}\RG{7.879}{-0.004}{10.9091}{rm}{91}\RG{10.910}{-0.004}{10.9091}{mi}{120}\RG{17.144}{-0.004}{10.9091}{rm}{93}\RG{20.175}{-1.640}{7.9701}{sy}{20}\RG{26.761}{-1.640}{7.9701}{mi}{110}%
}\hskip.1em\allowbreak{} 上定义

\[
\RB{173.846}{10.306}{3.225}{%
\RG{69.918}{4.506}{7.9701}{rm}{50}\RG{80.138}{4.506}{7.9701}{sy}{48}\RG{0.000}{0.002}{10.9091}{mi}{84}\RG{7.890}{0.002}{10.9091}{rm}{40}\RG{12.132}{0.002}{10.9091}{mi}{112}\RG{17.621}{0.002}{10.9091}{rm}{41}\RG{24.876}{0.002}{10.9091}{rm}{61}\RG{36.385}{0.002}{10.9091}{rm}{40}\RG{40.651}{0.002}{10.9091}{mi}{120}\RG{49.310}{0.002}{10.9091}{sy}{0}\RG{60.220}{0.002}{10.9091}{rm}{49}\RG{65.674}{0.002}{10.9091}{rm}{41}\RG{74.650}{0.002}{10.9091}{mi}{112}\RG{82.934}{0.002}{10.9091}{rm}{40}\RG{87.176}{0.002}{10.9091}{mi}{120}\RG{93.411}{0.002}{10.9091}{rm}{41}\RG{100.077}{0.002}{10.9091}{sy}{0}\RG{110.987}{0.002}{10.9091}{mi}{110}\RG{117.535}{0.002}{10.9091}{rm}{40}\RG{121.777}{0.002}{10.9091}{mi}{120}\RG{130.436}{0.002}{10.9091}{sy}{0}\RG{141.345}{0.002}{10.9091}{rm}{49}\RG{146.800}{0.002}{10.9091}{rm}{41}\RG{151.042}{0.002}{10.9091}{mi}{112}\RG{156.531}{0.002}{10.9091}{rm}{40}\RG{160.773}{0.002}{10.9091}{mi}{120}\RG{167.008}{0.002}{10.9091}{rm}{41}\RG{171.251}{0.002}{10.9091}{mi}{58}%
}
\]
对\allowbreak{}\hskip.1em \RB{46.391}{7.764}{2.738}{%
\RG{0.000}{-0.002}{10.9091}{rm}{48}\RG{8.485}{-0.002}{10.9091}{sy}{20}\RG{20.000}{-0.002}{10.9091}{mi}{106}\RG{28.147}{-0.002}{10.9091}{sy}{20}\RG{39.662}{-0.002}{10.9091}{mi}{110}%
}\hskip.1em\allowbreak{} ，有

\[
\RB{201.950}{10.307}{3.224}{%
\RG{18.367}{4.507}{7.9701}{mi}{106}\RG{52.889}{4.507}{7.9701}{mi}{106}\RG{56.773}{4.507}{7.9701}{sy}{0}\RG{63.360}{4.507}{7.9701}{rm}{49}\RG{126.599}{4.507}{7.9701}{mi}{106}\RG{184.033}{4.507}{7.9701}{mi}{106}\RG{187.917}{4.507}{7.9701}{rm}{43}\RG{194.500}{4.507}{7.9701}{rm}{49}\RG{0.000}{0.003}{10.9091}{mi}{84}\RG{7.890}{0.003}{10.9091}{rm}{40}\RG{12.132}{0.003}{10.9091}{mi}{120}\RG{22.749}{0.003}{10.9091}{rm}{41}\RG{30.004}{0.003}{10.9091}{rm}{61}\RG{41.537}{0.003}{10.9091}{mi}{106}\RG{46.643}{0.003}{10.9091}{mi}{120}\RG{70.517}{0.003}{10.9091}{rm}{43}\RG{81.415}{0.003}{10.9091}{rm}{40}\RG{85.668}{0.003}{10.9091}{mi}{110}\RG{94.641}{0.003}{10.9091}{sy}{0}\RG{105.550}{0.003}{10.9091}{rm}{50}\RG{111.004}{0.003}{10.9091}{mi}{106}\RG{116.121}{0.003}{10.9091}{rm}{41}\RG{120.364}{0.003}{10.9091}{mi}{120}\RG{133.405}{0.003}{10.9091}{rm}{43}\RG{144.303}{0.003}{10.9091}{rm}{40}\RG{148.556}{0.003}{10.9091}{mi}{106}\RG{156.097}{0.003}{10.9091}{sy}{0}\RG{167.007}{0.003}{10.9091}{mi}{110}\RG{173.555}{0.003}{10.9091}{rm}{41}\RG{177.797}{0.003}{10.9091}{mi}{120}\RG{199.236}{0.003}{10.9091}{mi}{59}%
}
\]
两端越界的项因系数为零而消失。因此\allowbreak{}\hskip.1em \RB{8.180}{7.888}{0.500}{%
\RG{0.000}{0.002}{10.9091}{mi}{84}%
}\hskip.1em\allowbreak{} 在基\allowbreak{}\hskip.1em \RB{52.393}{7.976}{2.603}{%
\RG{0.000}{0.002}{10.9091}{rm}{49}\RG{5.455}{0.002}{10.9091}{mi}{59}\RG{10.288}{0.002}{10.9091}{mi}{120}\RG{16.517}{0.002}{10.9091}{mi}{59}\RG{21.361}{0.002}{10.9091}{mi}{58}\RG{26.193}{0.002}{10.9091}{mi}{58}\RG{31.037}{0.002}{10.9091}{mi}{58}\RG{35.881}{0.002}{10.9091}{mi}{59}\RG{40.714}{0.002}{10.9091}{mi}{120}\RG{47.015}{3.961}{7.9701}{mi}{110}%
}\hskip.1em\allowbreak{} 下的矩阵恰为\allowbreak{}\hskip.1em \RB{14.158}{9.937}{0.500}{%
\RG{0.000}{0.002}{10.9091}{mi}{65}\RG{8.182}{3.961}{7.9701}{rm}{84}%
}\hskip.1em\allowbreak{} 。

改取基\allowbreak{}\hskip.1em \RB{98.455}{8.680}{3.229}{%
\RG{0.000}{-0.002}{10.9091}{rm}{49}\RG{5.455}{-0.002}{10.9091}{mi}{59}\RG{10.288}{-0.002}{10.9091}{mi}{120}\RG{18.962}{-0.002}{10.9091}{sy}{0}\RG{29.871}{-0.002}{10.9091}{rm}{49}\RG{35.326}{-0.002}{10.9091}{mi}{59}\RG{40.159}{-0.002}{10.9091}{mi}{58}\RG{45.003}{-0.002}{10.9091}{mi}{58}\RG{49.836}{-0.002}{10.9091}{mi}{58}\RG{54.679}{-0.002}{10.9091}{mi}{59}\RG{59.568}{-0.002}{10.9091}{rm}{40}\RG{63.812}{-0.002}{10.9091}{mi}{120}\RG{72.471}{-0.002}{10.9091}{sy}{0}\RG{83.380}{-0.002}{10.9091}{rm}{49}\RG{88.834}{-0.002}{10.9091}{rm}{41}\RG{93.077}{3.957}{7.9701}{mi}{110}%
}\hskip.1em\allowbreak{} ，则

\[
\RB{324.421}{10.299}{4.313}{%
\RG{7.890}{8.831}{10.9092}{ex}{0}\RG{50.780}{8.831}{10.9092}{ex}{1}\RG{252.682}{8.831}{10.9092}{ex}{0}\RG{296.826}{8.831}{10.9092}{ex}{1}\RG{46.398}{4.499}{7.9701}{mi}{106}\RG{137.317}{4.499}{7.9701}{mi}{106}\RG{141.201}{4.499}{7.9701}{rm}{43}\RG{147.784}{4.499}{7.9701}{rm}{49}\RG{291.189}{4.499}{7.9701}{mi}{110}\RG{0.000}{-0.005}{10.9091}{mi}{84}\RG{12.890}{-0.005}{10.9091}{rm}{40}\RG{17.133}{-0.005}{10.9091}{mi}{120}\RG{25.792}{-0.005}{10.9091}{sy}{0}\RG{36.701}{-0.005}{10.9091}{rm}{49}\RG{42.156}{-0.005}{10.9091}{rm}{41}\RG{58.810}{-0.005}{10.9091}{rm}{61}\RG{70.309}{-0.005}{10.9091}{rm}{40}\RG{74.569}{-0.005}{10.9091}{mi}{106}\RG{82.110}{-0.005}{10.9091}{sy}{0}\RG{93.019}{-0.005}{10.9091}{mi}{110}\RG{99.567}{-0.005}{10.9091}{rm}{41}\RG{103.800}{-0.005}{10.9091}{rm}{40}\RG{108.052}{-0.005}{10.9091}{mi}{120}\RG{116.711}{-0.005}{10.9091}{sy}{0}\RG{127.620}{-0.005}{10.9091}{rm}{49}\RG{133.075}{-0.005}{10.9091}{rm}{41}\RG{163.429}{-0.005}{10.9091}{rm}{40}\RG{167.662}{-0.005}{10.9091}{rm}{48}\RG{176.156}{-0.005}{10.9091}{sy}{20}\RG{187.672}{-0.005}{10.9091}{mi}{106}\RG{195.810}{-0.005}{10.9091}{mi}{60}\RG{207.308}{-0.005}{10.9091}{mi}{110}\RG{213.882}{-0.005}{10.9091}{rm}{41}\RG{218.125}{-0.005}{10.9091}{mi}{59}\RG{244.776}{-0.005}{10.9091}{mi}{84}\RG{257.682}{-0.005}{10.9091}{rm}{40}\RG{261.924}{-0.005}{10.9091}{mi}{120}\RG{270.583}{-0.005}{10.9091}{sy}{0}\RG{281.492}{-0.005}{10.9091}{rm}{49}\RG{286.947}{-0.005}{10.9091}{rm}{41}\RG{304.856}{-0.005}{10.9091}{rm}{61}\RG{316.355}{-0.005}{10.9091}{rm}{48}\RG{321.827}{-0.005}{10.9091}{mi}{58}%
}
\]
故\allowbreak{}\hskip.1em \RB{8.180}{7.887}{0.500}{%
\RG{0.000}{0.001}{10.9091}{mi}{84}%
}\hskip.1em\allowbreak{} 在这组基下的矩阵为严格下三角阵。于是

\[
\RB{203.893}{10.477}{3.230}{%
\RG{154.218}{4.501}{7.9701}{rm}{84}\RG{184.842}{4.501}{7.9701}{mi}{110}\RG{189.980}{4.501}{7.9701}{rm}{43}\RG{196.563}{4.501}{7.9701}{rm}{49}\RG{0.000}{-0.003}{10.9091}{rm}{100}\RG{6.055}{-0.003}{10.9091}{rm}{101}\RG{10.898}{-0.003}{10.9091}{rm}{116}\RG{15.131}{-0.003}{10.9091}{rm}{40}\RG{19.394}{-0.003}{10.9091}{mi}{122}\RG{24.936}{-0.003}{10.9091}{mi}{73}\RG{48.624}{-0.003}{10.9091}{sy}{0}\RG{59.533}{-0.003}{10.9091}{mi}{65}\RG{67.716}{-0.003}{10.9091}{rm}{41}\RG{74.970}{-0.003}{10.9091}{rm}{61}\RG{86.479}{-0.003}{10.9091}{rm}{100}\RG{92.534}{-0.003}{10.9091}{rm}{101}\RG{97.378}{-0.003}{10.9091}{rm}{116}\RG{101.610}{-0.003}{10.9091}{rm}{40}\RG{105.898}{-0.003}{10.9091}{mi}{122}\RG{111.440}{-0.003}{10.9091}{mi}{73}\RG{135.127}{-0.003}{10.9091}{sy}{0}\RG{146.036}{-0.003}{10.9091}{mi}{65}\RG{160.500}{-0.003}{10.9091}{rm}{41}\RG{167.755}{-0.003}{10.9091}{rm}{61}\RG{179.288}{-0.003}{10.9091}{mi}{122}\RG{201.299}{-0.003}{10.9091}{mi}{58}\RG{29.743}{-1.640}{7.9701}{mi}{110}\RG{34.881}{-1.640}{7.9701}{rm}{43}\RG{41.465}{-1.640}{7.9701}{rm}{49}\RG{116.246}{-1.640}{7.9701}{mi}{110}\RG{121.384}{-1.640}{7.9701}{rm}{43}\RG{127.967}{-1.640}{7.9701}{rm}{49}%
}
\]
因此\allowbreak{}\hskip.1em \RB{8.366}{8.314}{0.500}{%
\RG{0.000}{0.003}{10.9091}{mi}{65}%
}\hskip.1em\allowbreak{} 的唯一特征值是\allowbreak{}\hskip.1em \RB{5.518}{7.769}{0.737}{%
\RG{0.000}{0.003}{10.9091}{rm}{48}%
}\hskip.1em\allowbreak{} ，代数重数为\allowbreak{}\hskip.1em \RB{24.941}{7.769}{1.402}{%
\RG{0.000}{0.003}{10.9091}{mi}{110}\RG{8.972}{0.003}{10.9091}{rm}{43}\RG{19.870}{0.003}{10.9091}{rm}{49}%
}\hskip.1em\allowbreak{} 。

设\allowbreak{}\hskip.1em \RB{98.235}{9.934}{3.228}{%
\RG{0.000}{-0.001}{10.9091}{mi}{117}\RG{9.304}{-0.001}{10.9091}{rm}{61}\RG{20.835}{-0.001}{10.9091}{rm}{40}\RG{25.090}{-0.001}{10.9091}{mi}{117}\RG{31.335}{-1.637}{7.9701}{rm}{48}\RG{36.068}{-0.001}{10.9091}{mi}{59}\RG{40.901}{-0.001}{10.9091}{mi}{117}\RG{47.161}{-1.637}{7.9701}{rm}{49}\RG{51.893}{-0.001}{10.9091}{mi}{59}\RG{56.726}{-0.001}{10.9091}{mi}{58}\RG{61.570}{-0.001}{10.9091}{mi}{58}\RG{66.414}{-0.001}{10.9091}{mi}{58}\RG{71.246}{-0.001}{10.9091}{mi}{59}\RG{76.090}{-0.001}{10.9091}{mi}{117}\RG{82.382}{-1.637}{7.9701}{mi}{110}\RG{88.018}{-0.001}{10.9091}{rm}{41}\RG{92.260}{3.958}{7.9701}{rm}{84}%
}\hskip.1em\allowbreak{} 满足\allowbreak{}\hskip.1em \RB{34.535}{8.310}{0.741}{%
\RG{0.000}{-0.001}{10.9091}{mi}{65}\RG{8.182}{-0.001}{10.9091}{mi}{117}\RG{17.486}{-0.001}{10.9091}{rm}{61}\RG{29.017}{-0.001}{10.9091}{rm}{48}%
}\hskip.1em\allowbreak{} 。当\allowbreak{}\hskip.1em \RB{26.221}{7.765}{1.996}{%
\RG{0.000}{-0.001}{10.9091}{mi}{110}\RG{9.606}{-0.001}{10.9091}{sy}{21}\RG{21.150}{-0.001}{10.9091}{rm}{49}%
}\hskip.1em\allowbreak{} 时，第一行给出\allowbreak{}\hskip.1em \RB{35.888}{5.321}{2.304}{%
\RG{0.000}{-0.001}{10.9091}{mi}{117}\RG{6.245}{-1.637}{7.9701}{rm}{48}\RG{14.036}{-0.001}{10.9091}{rm}{61}\RG{25.580}{-0.001}{10.9091}{mi}{117}\RG{31.825}{-1.637}{7.9701}{rm}{49}%
}\hskip.1em\allowbreak{} 。若\allowbreak{}\hskip.1em \RB{46.087}{5.321}{3.763}{%
\RG{0.000}{-0.001}{10.9091}{mi}{117}\RG{6.245}{-1.637}{7.9701}{mi}{106}\RG{10.129}{-1.637}{7.9701}{sy}{0}\RG{16.715}{-1.637}{7.9701}{rm}{49}\RG{24.506}{-0.001}{10.9091}{rm}{61}\RG{36.050}{-0.001}{10.9091}{mi}{117}\RG{42.295}{-1.637}{7.9701}{mi}{106}%
}\hskip.1em\allowbreak{} ，则第\allowbreak{}\hskip.1em \RB{4.831}{7.715}{2.732}{%
\RG{0.000}{0.004}{10.9091}{mi}{106}%
}\hskip.1em\allowbreak{} 行化为

\[
\RB{186.605}{8.681}{3.763}{%
\RG{0.000}{-0.001}{10.9091}{rm}{40}\RG{4.242}{-0.001}{10.9091}{mi}{110}\RG{13.215}{-0.001}{10.9091}{sy}{0}\RG{24.124}{-0.001}{10.9091}{mi}{106}\RG{29.241}{-0.001}{10.9091}{rm}{41}\RG{33.474}{-0.001}{10.9091}{rm}{40}\RG{37.725}{-0.001}{10.9091}{mi}{117}\RG{50.777}{-0.001}{10.9091}{sy}{0}\RG{61.687}{-0.001}{10.9091}{mi}{117}\RG{83.135}{-0.001}{10.9091}{rm}{41}\RG{90.389}{-0.001}{10.9091}{rm}{61}\RG{101.898}{-0.001}{10.9091}{rm}{48}\RG{129.171}{-0.001}{10.9091}{rm}{40}\RG{133.404}{-0.001}{10.9091}{rm}{49}\RG{141.923}{-0.001}{10.9091}{sy}{20}\RG{153.438}{-0.001}{10.9091}{mi}{106}\RG{161.566}{-0.001}{10.9091}{mi}{60}\RG{173.075}{-0.001}{10.9091}{mi}{110}\RG{179.649}{-0.001}{10.9091}{rm}{41}\RG{183.891}{-0.001}{10.9091}{mi}{59}\RG{43.971}{-1.637}{7.9701}{mi}{106}\RG{67.932}{-1.637}{7.9701}{mi}{106}\RG{71.816}{-1.637}{7.9701}{rm}{43}\RG{78.399}{-1.637}{7.9701}{rm}{49}%
}
\]
从而依次得到\allowbreak{}\hskip.1em \RB{89.941}{5.320}{2.305}{%
\RG{0.000}{-0.002}{10.9091}{mi}{117}\RG{6.245}{-1.638}{7.9701}{rm}{48}\RG{14.008}{-0.002}{10.9091}{rm}{61}\RG{25.523}{-0.002}{10.9091}{mi}{117}\RG{31.768}{-1.638}{7.9701}{rm}{49}\RG{39.530}{-0.002}{10.9091}{rm}{61}\RG{51.045}{-0.002}{10.9091}{sy}{1}\RG{55.878}{-0.002}{10.9091}{sy}{1}\RG{60.722}{-0.002}{10.9091}{sy}{1}\RG{66.804}{-0.002}{10.9091}{rm}{61}\RG{78.319}{-0.002}{10.9091}{mi}{117}\RG{84.564}{-1.638}{7.9701}{mi}{110}%
}\hskip.1em\allowbreak{} 。反之，\allowbreak{}\hskip.1em \RB{8.366}{8.309}{0.502}{%
\RG{0.000}{-0.002}{10.9091}{mi}{65}%
}\hskip.1em\allowbreak{} 的每行元素之和均为零，故这些向量都在核中。因此全部特征向量为

\[
\RB{167.595}{10.478}{3.229}{%
\RG{79.452}{4.502}{7.9701}{rm}{84}\RG{0.000}{-0.002}{10.9091}{mi}{117}\RG{9.275}{-0.002}{10.9091}{rm}{61}\RG{20.791}{-0.002}{10.9091}{mi}{99}\RG{25.512}{-0.002}{10.9091}{rm}{40}\RG{29.745}{-0.002}{10.9091}{rm}{49}\RG{35.209}{-0.002}{10.9091}{mi}{59}\RG{40.057}{-0.002}{10.9091}{rm}{49}\RG{45.512}{-0.002}{10.9091}{mi}{59}\RG{50.345}{-0.002}{10.9091}{mi}{58}\RG{55.189}{-0.002}{10.9091}{mi}{58}\RG{60.022}{-0.002}{10.9091}{mi}{58}\RG{64.865}{-0.002}{10.9091}{mi}{59}\RG{69.755}{-0.002}{10.9091}{rm}{49}\RG{75.209}{-0.002}{10.9091}{rm}{41}\RG{85.734}{-0.002}{10.9091}{mi}{59}\RG{112.385}{-0.002}{10.9091}{mi}{99}\RG{120.152}{-0.002}{10.9091}{sy}{50}\RG{130.455}{-0.002}{10.9091}{bb}{67}\RG{140.758}{-0.002}{10.9091}{sy}{110}\RG{148.635}{-0.002}{10.9091}{sy}{102}\RG{154.091}{-0.002}{10.9091}{rm}{48}\RG{159.546}{-0.002}{10.9091}{sy}{103}\RG{165.000}{-0.002}{10.9091}{mi}{58}%
}
\]
当\allowbreak{}\hskip.1em \RB{26.595}{7.763}{0.743}{%
\RG{0.000}{-0.003}{10.9091}{mi}{110}\RG{9.578}{-0.003}{10.9091}{rm}{61}\RG{21.076}{-0.003}{10.9091}{rm}{48}%
}\hskip.1em\allowbreak{} 时，\allowbreak{}\hskip.1em \RB{36.050}{8.679}{3.230}{%
\RG{0.000}{-0.003}{10.9091}{mi}{65}\RG{11.212}{-0.003}{10.9091}{rm}{61}\RG{22.710}{-0.003}{10.9091}{rm}{40}\RG{26.943}{-0.003}{10.9091}{rm}{48}\RG{32.398}{-0.003}{10.9091}{rm}{41}%
}\hskip.1em\allowbreak{} ，上述结论仍成立。\allowbreak{}\hskip.1em \RB{8.366}{7.992}{0.503}{%
\RG{0.000}{-0.003}{10.9091}{am}{3}%
}\hskip.1em\allowbreak{}

方法小结：把带有等差系数的三对角阵识别为多项式算子，再平移多项式基，可直接看出幂零结构。

题目来源：谢启鸿《复旦大学高等代数习题集（2023版）》，第115页。

% Original page 2.
\section*{2022S03}
设\allowbreak{}\hskip.1em \RB{6.729}{5.319}{0.623}{%
\RG{0.000}{-0.003}{10.9091}{mi}{110}%
}\hskip.1em\allowbreak{} 阶复方阵\allowbreak{}\hskip.1em \RB{21.762}{8.308}{2.608}{%
\RG{0.000}{-0.003}{10.9091}{mi}{65}\RG{8.182}{-0.003}{10.9091}{mi}{59}\RG{13.015}{-0.003}{10.9091}{mi}{66}%
}\hskip.1em\allowbreak{} 满足\allowbreak{}\hskip.1em \RB{62.671}{8.308}{1.408}{%
\RG{0.000}{-0.003}{10.9091}{mi}{65}\RG{10.728}{-0.003}{10.9091}{rm}{43}\RG{21.760}{-0.003}{10.9091}{mi}{66}\RG{33.919}{-0.003}{10.9091}{rm}{61}\RG{45.741}{-0.003}{10.9091}{mi}{65}\RG{53.923}{-0.003}{10.9091}{mi}{66}%
}\hskip.1em\allowbreak{} 。证明：\allowbreak{}\hskip.1em \RB{8.366}{8.308}{0.503}{%
\RG{0.000}{-0.003}{10.9091}{mi}{65}%
}\hskip.1em\allowbreak{} 的特征值\allowbreak{}\hskip.1em \RB{47.080}{8.068}{2.608}{%
\RG{0.000}{-0.003}{10.9091}{mi}{21}\RG{6.364}{-1.639}{7.9701}{rm}{49}\RG{11.096}{-0.003}{10.9091}{mi}{59}\RG{15.929}{-0.003}{10.9091}{mi}{58}\RG{20.773}{-0.003}{10.9091}{mi}{58}\RG{25.616}{-0.003}{10.9091}{mi}{58}\RG{30.449}{-0.003}{10.9091}{mi}{59}\RG{35.293}{-0.003}{10.9091}{mi}{21}\RG{41.702}{-1.639}{7.9701}{mi}{110}%
}\hskip.1em\allowbreak{} 和\allowbreak{}\hskip.1em \RB{8.747}{7.948}{0.503}{%
\RG{0.000}{-0.003}{10.9091}{mi}{66}%
}\hskip.1em\allowbreak{} 的特征值\allowbreak{}\hskip.1em \RB{47.499}{5.319}{2.859}{%
\RG{0.000}{-0.003}{10.9091}{mi}{22}\RG{6.574}{-1.639}{7.9701}{rm}{49}\RG{11.306}{-0.003}{10.9091}{mi}{59}\RG{16.139}{-0.003}{10.9091}{mi}{58}\RG{20.983}{-0.003}{10.9091}{mi}{58}\RG{25.815}{-0.003}{10.9091}{mi}{58}\RG{30.659}{-0.003}{10.9091}{mi}{59}\RG{35.503}{-0.003}{10.9091}{mi}{22}\RG{42.121}{-1.639}{7.9701}{mi}{110}%
}\hskip.1em\allowbreak{}经过适当排序后，满足

\[
\RB{145.028}{8.679}{3.230}{%
\RG{0.000}{-0.003}{10.9091}{mi}{21}\RG{12.169}{-0.003}{10.9091}{rm}{43}\RG{23.078}{-0.003}{10.9091}{mi}{22}\RG{36.063}{-0.003}{10.9091}{rm}{61}\RG{47.578}{-0.003}{10.9091}{mi}{21}\RG{57.323}{-0.003}{10.9091}{mi}{22}\RG{89.096}{-0.003}{10.9091}{rm}{40}\RG{93.329}{-0.003}{10.9091}{rm}{49}\RG{101.824}{-0.003}{10.9091}{sy}{20}\RG{113.339}{-0.003}{10.9091}{mi}{105}\RG{120.127}{-0.003}{10.9091}{sy}{20}\RG{131.642}{-0.003}{10.9091}{mi}{110}\RG{138.190}{-0.003}{10.9091}{rm}{41}\RG{142.433}{-0.003}{10.9091}{mi}{58}\RG{6.363}{-1.639}{7.9701}{mi}{105}\RG{29.651}{-1.639}{7.9701}{mi}{105}\RG{53.941}{-1.639}{7.9701}{mi}{105}\RG{63.897}{-1.639}{7.9701}{mi}{105}%
}
\]
特别地，\allowbreak{}\hskip.1em \RB{8.366}{8.308}{0.503}{%
\RG{0.000}{-0.003}{10.9091}{mi}{65}%
}\hskip.1em\allowbreak{} 是幂零阵当且仅当\allowbreak{}\hskip.1em \RB{8.747}{7.948}{0.503}{%
\RG{0.000}{-0.003}{10.9091}{mi}{66}%
}\hskip.1em\allowbreak{} 是幂零阵。

证明：由题设有

\[
\RB{94.723}{8.680}{3.229}{%
\RG{0.000}{-0.002}{10.9091}{rm}{40}\RG{4.242}{-0.002}{10.9091}{mi}{73}\RG{12.318}{-0.002}{10.9091}{sy}{0}\RG{23.227}{-0.002}{10.9091}{mi}{65}\RG{31.409}{-0.002}{10.9091}{rm}{41}\RG{35.642}{-0.002}{10.9091}{rm}{40}\RG{39.894}{-0.002}{10.9091}{mi}{73}\RG{47.969}{-0.002}{10.9091}{sy}{0}\RG{58.878}{-0.002}{10.9091}{mi}{66}\RG{67.701}{-0.002}{10.9091}{rm}{41}\RG{74.955}{-0.002}{10.9091}{rm}{61}\RG{86.489}{-0.002}{10.9091}{mi}{73}\RG{92.129}{-0.002}{10.9091}{mi}{58}%
}
\]
取行列式可知\allowbreak{}\hskip.1em \RB{27.351}{8.308}{0.503}{%
\RG{0.000}{-0.003}{10.9091}{mi}{73}\RG{8.076}{-0.003}{10.9091}{sy}{0}\RG{18.985}{-0.003}{10.9091}{mi}{65}%
}\hskip.1em\allowbreak{} 与\allowbreak{}\hskip.1em \RB{27.733}{7.948}{0.503}{%
\RG{0.000}{-0.003}{10.9091}{mi}{73}\RG{8.076}{-0.003}{10.9091}{sy}{0}\RG{18.985}{-0.003}{10.9091}{mi}{66}%
}\hskip.1em\allowbreak{} 均可逆，故

\[
\RB{161.616}{10.307}{3.224}{%
\RG{78.004}{4.507}{7.9701}{sy}{0}\RG{84.591}{4.507}{7.9701}{rm}{49}\RG{147.702}{4.507}{7.9701}{sy}{0}\RG{154.288}{4.507}{7.9701}{rm}{49}\RG{0.000}{0.003}{10.9091}{mi}{66}\RG{11.852}{0.003}{10.9091}{rm}{61}\RG{23.367}{0.003}{10.9091}{mi}{73}\RG{31.443}{0.003}{10.9091}{sy}{0}\RG{42.352}{0.003}{10.9091}{rm}{40}\RG{46.594}{0.003}{10.9091}{mi}{73}\RG{54.670}{0.003}{10.9091}{sy}{0}\RG{65.580}{0.003}{10.9091}{mi}{65}\RG{73.762}{0.003}{10.9091}{rm}{41}\RG{92.353}{0.003}{10.9091}{rm}{61}\RG{103.869}{0.003}{10.9091}{mi}{65}\RG{112.050}{0.003}{10.9091}{rm}{40}\RG{116.293}{0.003}{10.9091}{mi}{65}\RG{126.899}{0.003}{10.9091}{sy}{0}\RG{137.808}{0.003}{10.9091}{mi}{73}\RG{143.459}{0.003}{10.9091}{rm}{41}\RG{159.021}{0.003}{10.9091}{mi}{58}%
}
\]
特别地，\allowbreak{}\hskip.1em \RB{5.071}{7.765}{0.501}{%
\RG{0.000}{-0.001}{10.9091}{rm}{49}%
}\hskip.1em\allowbreak{} 不是\allowbreak{}\hskip.1em \RB{8.366}{8.310}{0.501}{%
\RG{0.000}{-0.001}{10.9091}{mi}{65}%
}\hskip.1em\allowbreak{} 的特征值，且\allowbreak{}\hskip.1em \RB{8.366}{8.310}{0.501}{%
\RG{0.000}{-0.001}{10.9091}{mi}{65}%
}\hskip.1em\allowbreak{} 与\allowbreak{}\hskip.1em \RB{8.747}{7.950}{0.501}{%
\RG{0.000}{-0.001}{10.9091}{mi}{66}%
}\hskip.1em\allowbreak{} 可交换。

在复数域上，依次选取特征向量，并对商空间上的诱导变换作归纳，可以将任意方阵上三角化。故存在可逆矩阵\allowbreak{}\hskip.1em \RB{8.726}{7.951}{0.500}{%
\RG{0.000}{-0.000}{10.9091}{mi}{80}%
}\hskip.1em\allowbreak{} ，使

\[
\RB{61.217}{10.300}{2.610}{%
\RG{8.519}{4.500}{7.9701}{sy}{0}\RG{15.105}{4.500}{7.9701}{rm}{49}\RG{0.000}{-0.005}{10.9091}{mi}{80}\RG{19.838}{-0.005}{10.9091}{mi}{65}\RG{28.020}{-0.005}{10.9091}{mi}{80}\RG{39.569}{-0.005}{10.9091}{rm}{61}\RG{51.084}{-0.005}{10.9091}{mi}{85}\RG{58.503}{-0.005}{10.9091}{mi}{59}%
}
\]
其中\allowbreak{}\hskip.1em \RB{8.791}{7.954}{0.736}{%
\RG{0.000}{0.004}{10.9091}{mi}{85}%
}\hskip.1em\allowbreak{} 为上三角阵，主对角元按顺序记为\allowbreak{}\hskip.1em \RB{47.079}{8.074}{2.602}{%
\RG{0.000}{0.004}{10.9091}{mi}{21}\RG{6.363}{-1.633}{7.9701}{rm}{49}\RG{11.095}{0.004}{10.9091}{mi}{59}\RG{15.928}{0.004}{10.9091}{mi}{58}\RG{20.772}{0.004}{10.9091}{mi}{58}\RG{25.615}{0.004}{10.9091}{mi}{58}\RG{30.448}{0.004}{10.9091}{mi}{59}\RG{35.292}{0.004}{10.9091}{mi}{21}\RG{41.701}{-1.633}{7.9701}{mi}{110}%
}\hskip.1em\allowbreak{} 。于是

\[
\RB{117.467}{10.305}{3.226}{%
\RG{8.519}{4.505}{7.9701}{sy}{0}\RG{15.106}{4.505}{7.9701}{rm}{49}\RG{106.818}{4.505}{7.9701}{sy}{0}\RG{113.404}{4.505}{7.9701}{rm}{49}\RG{0.000}{0.001}{10.9091}{mi}{80}\RG{19.838}{0.001}{10.9091}{mi}{66}\RG{28.653}{0.001}{10.9091}{mi}{80}\RG{40.209}{0.001}{10.9091}{rm}{61}\RG{51.724}{0.001}{10.9091}{mi}{73}\RG{59.801}{0.001}{10.9091}{sy}{0}\RG{70.710}{0.001}{10.9091}{rm}{40}\RG{74.953}{0.001}{10.9091}{mi}{73}\RG{83.028}{0.001}{10.9091}{sy}{0}\RG{93.937}{0.001}{10.9091}{mi}{85}\RG{102.575}{0.001}{10.9091}{rm}{41}%
}
\]
也是上三角阵，其第\allowbreak{}\hskip.1em \RB{3.696}{7.708}{0.623}{%
\RG{0.000}{-0.003}{10.9091}{mi}{105}%
}\hskip.1em\allowbreak{} 个主对角元为

\[
\RB{122.276}{15.451}{9.699}{%
\RG{56.023}{7.380}{10.9091}{rm}{49}\RG{99.348}{7.380}{10.9091}{mi}{21}\RG{105.711}{5.744}{7.9701}{mi}{105}\RG{0.000}{0.000}{10.9091}{mi}{22}\RG{12.985}{0.000}{10.9091}{rm}{61}\RG{24.483}{0.000}{10.9091}{rm}{49}\RG{32.379}{0.000}{10.9091}{sy}{0}\RG{77.242}{0.000}{10.9091}{rm}{61}\RG{119.682}{0.000}{10.9091}{mi}{58}\RG{6.574}{-1.636}{7.9701}{mi}{105}\RG{44.484}{-7.483}{10.9091}{rm}{49}\RG{52.362}{-7.483}{10.9091}{sy}{0}\RG{63.271}{-7.483}{10.9091}{mi}{21}\RG{89.954}{-7.483}{10.9091}{mi}{21}\RG{102.123}{-7.483}{10.9091}{sy}{0}\RG{113.032}{-7.483}{10.9091}{rm}{49}\RG{69.635}{-9.119}{7.9701}{mi}{105}\RG{96.317}{-9.119}{7.9701}{mi}{105}\RR{44.484}{2.727}{28.532}{0.436}\RR{89.954}{2.727}{28.532}{0.436}%
}
\]
上三角阵的特征多项式是各对角因子之积，故这些\allowbreak{}\hskip.1em \RB{9.600}{5.325}{2.853}{%
\RG{0.000}{0.003}{10.9091}{mi}{22}\RG{6.573}{-1.633}{7.9701}{mi}{105}%
}\hskip.1em\allowbreak{} 正是\allowbreak{}\hskip.1em \RB{8.747}{7.954}{0.500}{%
\RG{0.000}{0.003}{10.9091}{mi}{66}%
}\hskip.1em\allowbreak{} 的全部特征值，按代数重数计。整理上式即得

\[
\RB{69.872}{8.074}{2.853}{%
\RG{0.000}{0.003}{10.9091}{mi}{21}\RG{12.169}{0.003}{10.9091}{rm}{43}\RG{23.078}{0.003}{10.9091}{mi}{22}\RG{36.063}{0.003}{10.9091}{rm}{61}\RG{47.578}{0.003}{10.9091}{mi}{21}\RG{57.323}{0.003}{10.9091}{mi}{22}\RG{67.278}{0.003}{10.9091}{mi}{58}\RG{6.364}{-1.634}{7.9701}{mi}{105}\RG{29.651}{-1.634}{7.9701}{mi}{105}\RG{53.942}{-1.634}{7.9701}{mi}{105}\RG{63.896}{-1.634}{7.9701}{mi}{105}%
}
\]
再证关于幂零性的结论。\allowbreak{}\hskip.1em \RB{8.366}{8.314}{0.500}{%
\RG{0.000}{0.003}{10.9091}{mi}{65}%
}\hskip.1em\allowbreak{} 与\allowbreak{}\hskip.1em \RB{46.301}{9.762}{3.224}{%
\RG{0.000}{0.003}{10.9091}{rm}{40}\RG{4.243}{0.003}{10.9091}{mi}{65}\RG{14.849}{0.003}{10.9091}{sy}{0}\RG{25.758}{0.003}{10.9091}{mi}{73}\RG{31.410}{0.003}{10.9091}{rm}{41}\RG{35.652}{3.962}{7.9701}{sy}{0}\RG{42.238}{3.962}{7.9701}{rm}{49}%
}\hskip.1em\allowbreak{} 可交换；若\allowbreak{}\hskip.1em \RB{32.783}{8.314}{0.737}{%
\RG{0.000}{0.003}{10.9091}{mi}{65}\RG{8.181}{3.962}{7.9701}{mi}{114}\RG{15.766}{0.003}{10.9091}{rm}{61}\RG{27.264}{0.003}{10.9091}{rm}{48}%
}\hskip.1em\allowbreak{} ，则

\[
\RB{110.046}{8.681}{3.228}{%
\RG{8.822}{4.504}{7.9701}{mi}{114}\RG{36.104}{4.504}{7.9701}{mi}{114}\RG{76.311}{4.504}{7.9701}{sy}{0}\RG{82.898}{4.504}{7.9701}{mi}{114}\RG{0.000}{-0.001}{10.9091}{mi}{66}\RG{16.407}{-0.001}{10.9091}{rm}{61}\RG{27.922}{-0.001}{10.9091}{mi}{65}\RG{40.659}{-0.001}{10.9091}{rm}{40}\RG{44.901}{-0.001}{10.9091}{mi}{65}\RG{55.507}{-0.001}{10.9091}{sy}{0}\RG{66.417}{-0.001}{10.9091}{mi}{73}\RG{72.069}{-0.001}{10.9091}{rm}{41}\RG{90.483}{-0.001}{10.9091}{rm}{61}\RG{101.981}{-0.001}{10.9091}{rm}{48}\RG{107.452}{-0.001}{10.9091}{mi}{58}%
}
\]
另一方面，由\allowbreak{}\hskip.1em \RB{92.412}{8.687}{3.222}{%
\RG{0.000}{0.005}{10.9091}{rm}{40}\RG{4.243}{0.005}{10.9091}{mi}{73}\RG{12.318}{0.005}{10.9091}{sy}{0}\RG{23.227}{0.005}{10.9091}{mi}{66}\RG{32.049}{0.005}{10.9091}{rm}{41}\RG{36.282}{0.005}{10.9091}{rm}{40}\RG{40.534}{0.005}{10.9091}{mi}{73}\RG{48.610}{0.005}{10.9091}{sy}{0}\RG{59.519}{0.005}{10.9091}{mi}{65}\RG{67.702}{0.005}{10.9091}{rm}{41}\RG{74.956}{0.005}{10.9091}{rm}{61}\RG{86.490}{0.005}{10.9091}{mi}{73}%
}\hskip.1em\allowbreak{} 同样可得\allowbreak{}\hskip.1em \RB{78.490}{9.764}{3.222}{%
\RG{0.000}{0.005}{10.9091}{mi}{65}\RG{11.212}{0.005}{10.9091}{rm}{61}\RG{22.727}{0.005}{10.9091}{mi}{66}\RG{31.549}{0.005}{10.9091}{rm}{40}\RG{35.792}{0.005}{10.9091}{mi}{66}\RG{47.038}{0.005}{10.9091}{sy}{0}\RG{57.947}{0.005}{10.9091}{mi}{73}\RG{63.599}{0.005}{10.9091}{rm}{41}\RG{67.842}{3.964}{7.9701}{sy}{0}\RG{74.428}{3.964}{7.9701}{rm}{49}%
}\hskip.1em\allowbreak{} ，所以\allowbreak{}\hskip.1em \RB{8.747}{7.956}{0.500}{%
\RG{0.000}{0.005}{10.9091}{mi}{66}%
}\hskip.1em\allowbreak{} 幂零也推出\allowbreak{}\hskip.1em \RB{8.366}{8.316}{0.500}{%
\RG{0.000}{0.005}{10.9091}{mi}{65}%
}\hskip.1em\allowbreak{} 幂零。\allowbreak{}\hskip.1em \RB{8.366}{8.000}{0.500}{%
\RG{0.000}{0.005}{10.9091}{am}{3}%
}\hskip.1em\allowbreak{}

方法小结：先把\allowbreak{}\hskip.1em \RB{62.511}{8.315}{1.401}{%
\RG{0.000}{0.004}{10.9091}{mi}{65}\RG{10.706}{0.004}{10.9091}{rm}{43}\RG{21.714}{0.004}{10.9091}{mi}{66}\RG{33.816}{0.004}{10.9091}{rm}{61}\RG{45.581}{0.004}{10.9091}{mi}{65}\RG{53.763}{0.004}{10.9091}{mi}{66}%
}\hskip.1em\allowbreak{} 改写为两个矩阵互为逆，再把\allowbreak{}\hskip.1em \RB{8.747}{7.955}{0.500}{%
\RG{0.000}{0.004}{10.9091}{mi}{66}%
}\hskip.1em\allowbreak{} 表成\allowbreak{}\hskip.1em \RB{8.366}{8.315}{0.500}{%
\RG{0.000}{0.004}{10.9091}{mi}{65}%
}\hskip.1em\allowbreak{} 的有理函数；特征值的配对随之确定。

题目来源：谢启鸿《复旦大学高等代数习题集（2023版）》，第115页。

% Original page 3.
\section*{2022S04}
设\allowbreak{}\hskip.1em \RB{55.861}{8.679}{3.230}{%
\RG{0.000}{-0.003}{10.9091}{mi}{65}\RG{11.212}{-0.003}{10.9091}{sy}{50}\RG{21.515}{-0.003}{10.9091}{mi}{77}\RG{32.098}{-1.639}{7.9701}{mi}{109}\RG{40.087}{-0.003}{10.9091}{rm}{40}\RG{44.329}{-0.003}{10.9091}{bb}{67}\RG{52.208}{-0.003}{10.9091}{rm}{41}%
}\hskip.1em\allowbreak{} 、\allowbreak{}\hskip.1em \RB{54.149}{8.679}{3.230}{%
\RG{0.000}{-0.003}{10.9091}{mi}{66}\RG{11.852}{-0.003}{10.9091}{sy}{50}\RG{22.155}{-0.003}{10.9091}{mi}{77}\RG{32.738}{-1.639}{7.9701}{mi}{110}\RG{38.374}{-0.003}{10.9091}{rm}{40}\RG{42.617}{-0.003}{10.9091}{bb}{67}\RG{50.496}{-0.003}{10.9091}{rm}{41}%
}\hskip.1em\allowbreak{} 的特征值分别为\allowbreak{}\hskip.1em \RB{49.431}{8.068}{2.608}{%
\RG{0.000}{-0.003}{10.9091}{mi}{21}\RG{6.364}{-1.639}{7.9701}{rm}{49}\RG{11.096}{-0.003}{10.9091}{mi}{59}\RG{15.929}{-0.003}{10.9091}{mi}{58}\RG{20.773}{-0.003}{10.9091}{mi}{58}\RG{25.616}{-0.003}{10.9091}{mi}{58}\RG{30.449}{-0.003}{10.9091}{mi}{59}\RG{35.293}{-0.003}{10.9091}{mi}{21}\RG{41.702}{-1.639}{7.9701}{mi}{109}%
}\hskip.1em\allowbreak{} 和\allowbreak{}\hskip.1em \RB{47.500}{5.319}{2.859}{%
\RG{0.000}{-0.003}{10.9091}{mi}{22}\RG{6.574}{-1.639}{7.9701}{rm}{49}\RG{11.306}{-0.003}{10.9091}{mi}{59}\RG{16.139}{-0.003}{10.9091}{mi}{58}\RG{20.983}{-0.003}{10.9091}{mi}{58}\RG{25.815}{-0.003}{10.9091}{mi}{58}\RG{30.659}{-0.003}{10.9091}{mi}{59}\RG{35.503}{-0.003}{10.9091}{mi}{22}\RG{42.122}{-1.639}{7.9701}{mi}{110}%
}\hskip.1em\allowbreak{} 。利用题设的可对角化摄动法，不用Kronecker积，求\allowbreak{}\hskip.1em \RB{69.404}{8.684}{3.225}{%
\RG{0.000}{0.002}{10.9091}{mi}{86}\RG{11.818}{0.002}{10.9091}{rm}{61}\RG{23.333}{0.002}{10.9091}{mi}{77}\RG{33.916}{-1.634}{7.9701}{mi}{109}\RG{41.407}{-1.634}{7.9701}{sy}{2}\RG{47.993}{-1.634}{7.9701}{mi}{110}\RG{53.630}{0.002}{10.9091}{rm}{40}\RG{57.872}{0.002}{10.9091}{bb}{67}\RG{65.752}{0.002}{10.9091}{rm}{41}%
}\hskip.1em\allowbreak{} 上下列线性变换的特征值：

\[
\RB{245.556}{9.291}{3.229}{%
\RT{0.000}{-0.002}{10.9091}{(}\RT{4.244}{-0.002}{10.9091}{1}\RT{9.698}{-0.002}{10.9091}{)}\RG{24.851}{-0.002}{10.9091}{mi}{39}\RG{36.720}{-0.002}{10.9091}{rm}{40}\RG{40.962}{-0.002}{10.9091}{mi}{88}\RG{50.856}{-0.002}{10.9091}{rm}{41}\RG{58.111}{-0.002}{10.9091}{rm}{61}\RG{69.645}{-0.002}{10.9091}{mi}{65}\RG{77.826}{-0.002}{10.9091}{mi}{88}\RG{87.710}{-0.002}{10.9091}{mi}{66}\RG{96.543}{-0.002}{10.9091}{rm}{59}\RT{123.210}{-0.002}{10.9091}{(}\RT{127.454}{-0.002}{10.9091}{2}\RT{132.908}{-0.002}{10.9091}{)}\RG{148.060}{-0.002}{10.9091}{mi}{39}\RG{159.929}{-0.002}{10.9091}{rm}{40}\RG{164.172}{-0.002}{10.9091}{mi}{88}\RG{174.065}{-0.002}{10.9091}{rm}{41}\RG{181.320}{-0.002}{10.9091}{rm}{61}\RG{192.854}{-0.002}{10.9091}{mi}{65}\RG{201.035}{-0.002}{10.9091}{mi}{88}\RG{213.354}{-0.002}{10.9091}{sy}{0}\RG{224.263}{-0.002}{10.9091}{mi}{88}\RG{234.147}{-0.002}{10.9091}{mi}{66}\RG{242.961}{-0.002}{10.9091}{mi}{58}\RG{31.988}{-1.638}{7.9701}{rm}{49}\RG{155.197}{-1.638}{7.9701}{rm}{50}%
}
\]
题设的摄动法是：先将矩阵上三角化，再给对角元加上\allowbreak{}\hskip.1em \RB{12.202}{7.333}{2.212}{%
\RG{0.000}{0.004}{10.9091}{mi}{99}\RG{4.721}{-1.632}{7.9701}{mi}{105}\RG{8.102}{0.004}{10.9091}{mi}{116}%
}\hskip.1em\allowbreak{} ，使小的正数\allowbreak{}\hskip.1em \RB{4.100}{7.333}{0.616}{%
\RG{0.000}{0.004}{10.9091}{mi}{116}%
}\hskip.1em\allowbreak{} 对应的矩阵具有互异特征值，并在\allowbreak{}\hskip.1em \RB{32.969}{9.272}{0.741}{%
\RG{0.000}{-0.001}{10.9091}{mi}{116}\RG{6.970}{-0.001}{10.9091}{sy}{33}\RG{20.909}{-0.001}{10.9091}{rm}{48}\RG{26.364}{3.958}{7.9701}{rm}{43}%
}\hskip.1em\allowbreak{} 时趋于原矩阵。解：取可逆矩阵\allowbreak{}\hskip.1em \RB{20.715}{8.187}{2.620}{%
\RG{0.000}{-0.004}{10.9091}{mi}{80}\RG{7.298}{-0.004}{10.9091}{mi}{59}\RG{12.142}{-0.004}{10.9091}{mi}{81}%
}\hskip.1em\allowbreak{} ，使\allowbreak{}\hskip.1em \RB{59.929}{9.754}{0.744}{%
\RG{0.000}{-0.004}{10.9091}{mi}{85}\RG{11.668}{-0.004}{10.9091}{rm}{61}\RG{23.183}{-0.004}{10.9091}{mi}{80}\RG{31.702}{3.954}{7.9701}{sy}{0}\RG{38.289}{3.954}{7.9701}{rm}{49}\RG{43.021}{-0.004}{10.9091}{mi}{65}\RG{51.203}{-0.004}{10.9091}{mi}{80}%
}\hskip.1em\allowbreak{} 、\allowbreak{}\hskip.1em \RB{63.695}{9.754}{2.620}{%
\RG{0.000}{-0.004}{10.9091}{mi}{87}\RG{14.849}{-0.004}{10.9091}{rm}{61}\RG{26.364}{-0.004}{10.9091}{mi}{81}\RG{34.988}{3.954}{7.9701}{sy}{0}\RG{41.574}{3.954}{7.9701}{rm}{49}\RG{46.307}{-0.004}{10.9091}{mi}{66}\RG{55.122}{-0.004}{10.9091}{mi}{81}%
}\hskip.1em\allowbreak{} 为上三角阵，其对角元按上述顺序排列。令

\[
\RB{350.278}{10.306}{4.306}{%
\RG{34.803}{8.838}{10.9092}{ex}{0}\RG{135.444}{8.838}{10.9092}{ex}{1}\RG{221.950}{8.838}{10.9092}{ex}{0}\RG{322.741}{8.838}{10.9092}{ex}{1}\RG{148.963}{4.506}{7.9701}{sy}{0}\RG{155.549}{4.506}{7.9701}{rm}{49}\RG{336.365}{4.506}{7.9701}{sy}{0}\RG{342.952}{4.506}{7.9701}{rm}{49}\RG{0.000}{0.002}{10.9091}{mi}{65}\RG{14.769}{0.002}{10.9091}{rm}{61}\RG{26.284}{0.002}{10.9091}{mi}{80}\RG{39.803}{0.002}{10.9091}{mi}{85}\RG{50.865}{0.002}{10.9091}{rm}{43}\RG{61.774}{0.002}{10.9091}{mi}{116}\RG{67.532}{0.002}{10.9091}{rm}{100}\RG{73.586}{0.002}{10.9091}{rm}{105}\RG{76.608}{0.002}{10.9091}{rm}{97}\RG{82.063}{0.002}{10.9091}{rm}{103}\RG{87.659}{0.002}{10.9091}{rm}{40}\RG{91.892}{0.002}{10.9091}{rm}{49}\RG{97.381}{0.002}{10.9091}{mi}{59}\RG{102.214}{0.002}{10.9091}{mi}{58}\RG{107.057}{0.002}{10.9091}{mi}{58}\RG{111.890}{0.002}{10.9091}{mi}{58}\RG{116.734}{0.002}{10.9091}{mi}{59}\RG{121.577}{0.002}{10.9091}{mi}{109}\RG{131.202}{0.002}{10.9091}{rm}{41}\RG{140.444}{0.002}{10.9091}{mi}{80}\RG{160.282}{0.002}{10.9091}{mi}{59}\RG{186.933}{0.002}{10.9091}{mi}{66}\RG{201.810}{0.002}{10.9091}{rm}{61}\RG{213.326}{0.002}{10.9091}{mi}{81}\RG{226.950}{0.002}{10.9091}{mi}{87}\RG{241.192}{0.002}{10.9091}{rm}{43}\RG{252.101}{0.002}{10.9091}{mi}{116}\RG{257.859}{0.002}{10.9091}{rm}{100}\RG{263.914}{0.002}{10.9091}{rm}{105}\RG{266.936}{0.002}{10.9091}{rm}{97}\RG{272.390}{0.002}{10.9091}{rm}{103}\RG{277.987}{0.002}{10.9091}{rm}{40}\RG{282.219}{0.002}{10.9091}{rm}{49}\RG{287.707}{0.002}{10.9091}{mi}{59}\RG{292.540}{0.002}{10.9091}{mi}{58}\RG{297.384}{0.002}{10.9091}{mi}{58}\RG{302.227}{0.002}{10.9091}{mi}{58}\RG{307.060}{0.002}{10.9091}{mi}{59}\RG{311.904}{0.002}{10.9091}{mi}{110}\RG{318.499}{0.002}{10.9091}{rm}{41}\RG{327.741}{0.002}{10.9091}{mi}{81}\RG{347.684}{0.002}{10.9091}{mi}{58}\RG{8.182}{-1.635}{7.9701}{mi}{116}\RG{195.224}{-1.635}{7.9701}{mi}{116}%
}
\]
当\allowbreak{}\hskip.1em \RB{24.288}{7.764}{0.938}{%
\RG{0.000}{-0.002}{10.9091}{mi}{116}\RG{7.102}{-0.002}{10.9091}{mi}{62}\RG{18.770}{-0.002}{10.9091}{rm}{48}%
}\hskip.1em\allowbreak{} 足够小时，\allowbreak{}\hskip.1em \RB{31.045}{8.069}{2.219}{%
\RG{0.000}{-0.002}{10.9091}{mi}{21}\RG{6.364}{-1.639}{7.9701}{mi}{105}\RG{12.226}{-0.002}{10.9091}{rm}{43}\RG{23.192}{-0.002}{10.9091}{mi}{105}\RG{26.945}{-0.002}{10.9091}{mi}{116}%
}\hskip.1em\allowbreak{} 两两不同，\allowbreak{}\hskip.1em \RB{33.609}{7.709}{3.765}{%
\RG{0.000}{-0.002}{10.9091}{mi}{22}\RG{6.573}{-1.639}{7.9701}{mi}{106}\RG{13.437}{-0.002}{10.9091}{rm}{43}\RG{24.403}{-0.002}{10.9091}{mi}{106}\RG{29.509}{-0.002}{10.9091}{mi}{116}%
}\hskip.1em\allowbreak{} 也两两不同：原来相等的一对被不同斜率分开，原来不同的一对则在充分小的邻域内保持不同。因此\allowbreak{}\hskip.1em \RB{28.151}{8.314}{2.602}{%
\RG{0.000}{0.003}{10.9091}{mi}{65}\RG{8.181}{-1.634}{7.9701}{mi}{116}\RG{11.738}{0.003}{10.9091}{mi}{59}\RG{16.571}{0.003}{10.9091}{mi}{66}\RG{24.861}{-1.634}{7.9701}{mi}{116}%
}\hskip.1em\allowbreak{} 均可对角化，且分别趋于\allowbreak{}\hskip.1em \RB{21.762}{8.314}{2.602}{%
\RG{0.000}{0.003}{10.9091}{mi}{65}\RG{8.182}{0.003}{10.9091}{mi}{59}\RG{13.015}{0.003}{10.9091}{mi}{66}%
}\hskip.1em\allowbreak{} 。

固定这样的\allowbreak{}\hskip.1em \RB{4.100}{7.328}{0.621}{%
\RG{0.000}{-0.001}{10.9091}{mi}{116}%
}\hskip.1em\allowbreak{} 。取\allowbreak{}\hskip.1em \RB{11.471}{8.310}{2.217}{%
\RG{0.000}{-0.001}{10.9091}{mi}{65}\RG{8.181}{-1.637}{7.9701}{mi}{116}%
}\hskip.1em\allowbreak{} 的特征向量基\allowbreak{}\hskip.1em \RB{48.236}{5.321}{2.606}{%
\RG{0.000}{-0.001}{10.9091}{mi}{97}\RG{5.766}{-1.637}{7.9701}{rm}{49}\RG{10.498}{-0.001}{10.9091}{mi}{59}\RG{15.331}{-0.001}{10.9091}{mi}{58}\RG{20.175}{-0.001}{10.9091}{mi}{58}\RG{25.018}{-0.001}{10.9091}{mi}{58}\RG{29.851}{-0.001}{10.9091}{mi}{59}\RG{34.695}{-0.001}{10.9091}{mi}{97}\RG{40.507}{-1.637}{7.9701}{mi}{109}%
}\hskip.1em\allowbreak{} ，以及\allowbreak{}\hskip.1em \RB{14.798}{9.934}{3.278}{%
\RG{0.000}{-0.001}{10.9091}{mi}{66}\RG{8.274}{-2.698}{7.9701}{mi}{116}\RG{8.822}{3.958}{7.9701}{rm}{84}%
}\hskip.1em\allowbreak{} 的特征向量基\allowbreak{}\hskip.1em \RB{43.716}{8.070}{2.606}{%
\RG{0.000}{-0.001}{10.9091}{mi}{98}\RG{4.682}{-1.637}{7.9701}{rm}{49}\RG{9.414}{-0.001}{10.9091}{mi}{59}\RG{14.247}{-0.001}{10.9091}{mi}{58}\RG{19.091}{-0.001}{10.9091}{mi}{58}\RG{23.923}{-0.001}{10.9091}{mi}{58}\RG{28.767}{-0.001}{10.9091}{mi}{59}\RG{33.611}{-0.001}{10.9091}{mi}{98}\RG{38.338}{-1.637}{7.9701}{mi}{110}%
}\hskip.1em\allowbreak{} ，使

\[
\RB{202.709}{10.475}{3.767}{%
\RG{119.329}{4.499}{7.9701}{rm}{84}\RG{0.000}{-0.005}{10.9091}{mi}{65}\RG{11.738}{-0.005}{10.9091}{mi}{97}\RG{23.916}{-0.005}{10.9091}{rm}{61}\RG{35.414}{-0.005}{10.9091}{rm}{40}\RG{39.673}{-0.005}{10.9091}{mi}{21}\RG{51.843}{-0.005}{10.9091}{rm}{43}\RG{62.753}{-0.005}{10.9091}{mi}{105}\RG{66.505}{-0.005}{10.9091}{mi}{116}\RG{70.450}{-0.005}{10.9091}{rm}{41}\RG{74.693}{-0.005}{10.9091}{mi}{97}\RG{83.841}{-0.005}{10.9091}{mi}{59}\RG{110.492}{-0.005}{10.9091}{mi}{66}\RG{125.611}{-0.005}{10.9091}{mi}{98}\RG{137.705}{-0.005}{10.9091}{rm}{61}\RG{149.203}{-0.005}{10.9091}{rm}{40}\RG{153.463}{-0.005}{10.9091}{mi}{22}\RG{166.843}{-0.005}{10.9091}{rm}{43}\RG{177.753}{-0.005}{10.9091}{mi}{106}\RG{182.858}{-0.005}{10.9091}{mi}{116}\RG{186.809}{-0.005}{10.9091}{rm}{41}\RG{191.051}{-0.005}{10.9091}{mi}{98}\RG{200.115}{-0.005}{10.9091}{mi}{58}\RG{8.182}{-1.641}{7.9701}{mi}{116}\RG{17.504}{-1.641}{7.9701}{mi}{105}\RG{46.037}{-1.641}{7.9701}{mi}{105}\RG{80.459}{-1.641}{7.9701}{mi}{105}\RG{130.293}{-1.641}{7.9701}{mi}{106}\RG{160.036}{-1.641}{7.9701}{mi}{106}\RG{195.733}{-1.641}{7.9701}{mi}{106}\RG{118.782}{-2.702}{7.9701}{mi}{116}%
}
\]
矩阵\allowbreak{}\hskip.1em \RB{51.155}{9.936}{5.180}{%
\RG{0.000}{0.001}{10.9091}{mi}{88}\RG{9.038}{-1.635}{7.9701}{mi}{105}\RG{11.915}{-1.635}{7.9701}{mi}{106}\RG{19.584}{0.001}{10.9091}{rm}{61}\RG{31.350}{0.001}{10.9091}{mi}{97}\RG{37.116}{-1.635}{7.9701}{mi}{105}\RG{40.497}{0.001}{10.9091}{mi}{98}\RG{45.179}{3.960}{7.9701}{rm}{84}\RG{45.179}{-3.054}{7.9701}{mi}{106}%
}\hskip.1em\allowbreak{} 构成\allowbreak{}\hskip.1em \RB{8.889}{7.952}{0.739}{%
\RG{0.000}{0.001}{10.9091}{mi}{86}%
}\hskip.1em\allowbreak{} 的一组基。事实上，若\allowbreak{}\hskip.1em \RB{79.806}{8.683}{3.226}{%
\RG{0.000}{0.001}{10.9091}{mi}{82}\RG{11.648}{0.001}{10.9091}{rm}{61}\RG{23.397}{0.001}{10.9091}{rm}{40}\RG{27.657}{0.001}{10.9091}{mi}{97}\RG{33.423}{-1.635}{7.9701}{rm}{49}\RG{38.156}{0.001}{10.9091}{mi}{59}\RG{42.989}{0.001}{10.9091}{mi}{58}\RG{47.833}{0.001}{10.9091}{mi}{58}\RG{52.666}{0.001}{10.9091}{mi}{58}\RG{57.509}{0.001}{10.9091}{mi}{59}\RG{62.353}{0.001}{10.9091}{mi}{97}\RG{68.165}{-1.635}{7.9701}{mi}{109}\RG{76.154}{0.001}{10.9091}{rm}{41}%
}\hskip.1em\allowbreak{} 、\allowbreak{}\hskip.1em \RB{74.236}{8.683}{3.226}{%
\RG{0.000}{0.001}{10.9091}{mi}{83}\RG{10.600}{0.001}{10.9091}{rm}{61}\RG{22.349}{0.001}{10.9091}{rm}{40}\RG{26.608}{0.001}{10.9091}{mi}{98}\RG{31.290}{-1.635}{7.9701}{rm}{49}\RG{36.022}{0.001}{10.9091}{mi}{59}\RG{40.855}{0.001}{10.9091}{mi}{58}\RG{45.699}{0.001}{10.9091}{mi}{58}\RG{50.532}{0.001}{10.9091}{mi}{58}\RG{55.375}{0.001}{10.9091}{mi}{59}\RG{60.219}{0.001}{10.9091}{mi}{98}\RG{64.947}{-1.635}{7.9701}{mi}{110}\RG{70.584}{0.001}{10.9091}{rm}{41}%
}\hskip.1em\allowbreak{} ，则可逆线性映射\allowbreak{}\hskip.1em \RB{56.146}{9.931}{0.744}{%
\RG{0.000}{-0.004}{10.9091}{mi}{89}\RG{11.788}{-0.004}{10.9091}{sy}{55}\RG{11.788}{-0.004}{10.9091}{sy}{33}\RG{25.727}{-0.004}{10.9091}{mi}{82}\RG{34.084}{-0.004}{10.9091}{mi}{89}\RG{42.833}{-0.004}{10.9091}{mi}{83}\RG{50.170}{3.955}{7.9701}{rm}{84}%
}\hskip.1em\allowbreak{} 将标准基\allowbreak{}\hskip.1em \RB{14.722}{7.914}{3.766}{%
\RG{0.000}{-0.004}{10.9091}{mi}{69}\RG{8.053}{-1.640}{7.9701}{mi}{105}\RG{10.930}{-1.640}{7.9701}{mi}{106}%
}\hskip.1em\allowbreak{} 映为\allowbreak{}\hskip.1em \RB{15.707}{7.947}{3.766}{%
\RG{0.000}{-0.004}{10.9091}{mi}{88}\RG{9.038}{-1.640}{7.9701}{mi}{105}\RG{11.915}{-1.640}{7.9701}{mi}{106}%
}\hskip.1em\allowbreak{} 。直接计算得

\[
\RB{194.953}{8.684}{22.469}{%
\RG{29.637}{0.002}{10.9091}{mi}{65}\RG{41.375}{0.002}{10.9091}{mi}{88}\RG{57.678}{0.002}{10.9091}{mi}{66}\RG{72.540}{0.002}{10.9091}{rm}{61}\RG{84.038}{0.002}{10.9091}{rm}{40}\RG{88.298}{0.002}{10.9091}{mi}{21}\RG{100.467}{0.002}{10.9091}{rm}{43}\RG{111.376}{0.002}{10.9091}{mi}{105}\RG{115.129}{0.002}{10.9091}{mi}{116}\RG{119.073}{0.002}{10.9091}{rm}{41}\RG{123.306}{0.002}{10.9091}{rm}{40}\RG{127.558}{0.002}{10.9091}{mi}{22}\RG{140.938}{0.002}{10.9091}{rm}{43}\RG{151.847}{0.002}{10.9091}{mi}{106}\RG{156.953}{0.002}{10.9091}{mi}{116}\RG{160.903}{0.002}{10.9091}{rm}{41}\RG{165.146}{0.002}{10.9091}{mi}{88}\RG{181.449}{0.002}{10.9091}{mi}{59}\RG{37.819}{-1.634}{7.9701}{mi}{116}\RG{50.413}{-1.634}{7.9701}{mi}{105}\RG{53.291}{-1.634}{7.9701}{mi}{106}\RG{65.954}{-1.634}{7.9701}{mi}{116}\RG{94.661}{-1.634}{7.9701}{mi}{105}\RG{134.132}{-1.634}{7.9701}{mi}{106}\RG{174.184}{-1.634}{7.9701}{mi}{105}\RG{177.062}{-1.634}{7.9701}{mi}{106}\RG{84.055}{-9.325}{10.9092}{ex}{0}\RG{171.055}{-9.325}{10.9092}{ex}{1}\RG{0.000}{-18.162}{10.9091}{mi}{65}\RG{11.738}{-18.162}{10.9091}{mi}{88}\RG{30.466}{-18.162}{10.9091}{sy}{0}\RG{41.375}{-18.162}{10.9091}{mi}{88}\RG{57.678}{-18.162}{10.9091}{mi}{66}\RG{72.540}{-18.162}{10.9091}{rm}{61}\RG{89.055}{-18.162}{10.9091}{mi}{21}\RG{101.224}{-18.162}{10.9091}{sy}{0}\RG{112.133}{-18.162}{10.9091}{mi}{22}\RG{125.513}{-18.162}{10.9091}{rm}{43}\RG{136.411}{-18.162}{10.9091}{rm}{40}\RG{140.664}{-18.162}{10.9091}{mi}{105}\RG{146.847}{-18.162}{10.9091}{sy}{0}\RG{157.756}{-18.162}{10.9091}{mi}{106}\RG{162.873}{-18.162}{10.9091}{rm}{41}\RG{167.115}{-18.162}{10.9091}{mi}{116}\RG{176.055}{-18.162}{10.9091}{mi}{88}\RG{192.359}{-18.162}{10.9091}{mi}{58}\RG{8.182}{-19.798}{7.9701}{mi}{116}\RG{20.776}{-19.798}{7.9701}{mi}{105}\RG{23.653}{-19.798}{7.9701}{mi}{106}\RG{50.413}{-19.798}{7.9701}{mi}{105}\RG{53.291}{-19.798}{7.9701}{mi}{106}\RG{65.954}{-19.798}{7.9701}{mi}{116}\RG{95.419}{-19.798}{7.9701}{mi}{105}\RG{118.707}{-19.798}{7.9701}{mi}{106}\RG{185.094}{-19.798}{7.9701}{mi}{105}\RG{187.971}{-19.798}{7.9701}{mi}{106}%
}
\]
所以，对\allowbreak{}\hskip.1em \RB{83.400}{8.686}{3.670}{%
\RG{0.000}{0.004}{10.9091}{mi}{39}\RG{7.136}{-1.632}{7.9701}{rm}{49}\RG{11.371}{-1.632}{7.9701}{mi}{59}\RG{13.722}{-1.632}{7.9701}{mi}{116}\RG{17.279}{0.004}{10.9091}{rm}{40}\RG{21.522}{0.004}{10.9091}{mi}{88}\RG{31.415}{0.004}{10.9091}{rm}{41}\RG{38.670}{0.004}{10.9091}{rm}{61}\RG{50.203}{0.004}{10.9091}{mi}{65}\RG{58.385}{-1.632}{7.9701}{mi}{116}\RG{61.942}{0.004}{10.9091}{mi}{88}\RG{71.825}{0.004}{10.9091}{mi}{66}\RG{80.111}{-1.632}{7.9701}{mi}{116}%
}\hskip.1em\allowbreak{} 、\allowbreak{}\hskip.1em \RB{106.627}{8.686}{3.670}{%
\RG{0.000}{0.004}{10.9091}{mi}{39}\RG{7.136}{-1.632}{7.9701}{rm}{50}\RG{11.371}{-1.632}{7.9701}{mi}{59}\RG{13.722}{-1.632}{7.9701}{mi}{116}\RG{17.279}{0.004}{10.9091}{rm}{40}\RG{21.522}{0.004}{10.9091}{mi}{88}\RG{31.416}{0.004}{10.9091}{rm}{41}\RG{38.671}{0.004}{10.9091}{rm}{61}\RG{50.203}{0.004}{10.9091}{mi}{65}\RG{58.386}{-1.632}{7.9701}{mi}{116}\RG{61.942}{0.004}{10.9091}{mi}{88}\RG{74.260}{0.004}{10.9091}{sy}{0}\RG{85.170}{0.004}{10.9091}{mi}{88}\RG{95.053}{0.004}{10.9091}{mi}{66}\RG{103.338}{-1.632}{7.9701}{mi}{116}%
}\hskip.1em\allowbreak{} ，有

\[
\RB{229.489}{17.656}{53.081}{%
\RG{89.680}{13.641}{7.9701}{mi}{109}\RG{106.996}{13.641}{7.9701}{mi}{110}\RG{86.455}{10.368}{10.9092}{ex}{89}\RG{102.595}{10.368}{10.9092}{ex}{89}\RG{116.917}{8.841}{10.9092}{ex}{2}\RG{221.441}{8.841}{10.9092}{ex}{3}\RG{0.000}{0.004}{10.9091}{rm}{100}\RG{6.055}{0.004}{10.9091}{rm}{101}\RG{10.898}{0.004}{10.9091}{rm}{116}\RG{15.131}{0.004}{10.9091}{rm}{40}\RG{19.394}{0.004}{10.9091}{mi}{122}\RG{24.936}{0.004}{10.9091}{mi}{73}\RG{39.478}{0.004}{10.9091}{sy}{0}\RG{50.387}{0.004}{10.9091}{mi}{39}\RG{67.667}{0.004}{10.9091}{rm}{41}\RG{74.921}{0.004}{10.9091}{rm}{61}\RG{121.463}{0.004}{10.9091}{mi}{122}\RG{129.440}{0.004}{10.9091}{sy}{0}\RG{140.349}{0.004}{10.9091}{rm}{40}\RG{144.592}{0.004}{10.9091}{mi}{21}\RG{156.761}{0.004}{10.9091}{rm}{43}\RG{167.670}{0.004}{10.9091}{mi}{105}\RG{171.423}{0.004}{10.9091}{mi}{116}\RG{175.368}{0.004}{10.9091}{rm}{41}\RG{179.601}{0.004}{10.9091}{rm}{40}\RG{183.853}{0.004}{10.9091}{mi}{22}\RG{197.233}{0.004}{10.9091}{rm}{43}\RG{208.142}{0.004}{10.9091}{mi}{106}\RG{213.247}{0.004}{10.9091}{mi}{116}\RG{217.199}{0.004}{10.9091}{rm}{41}\RG{225.986}{0.004}{10.9091}{mi}{59}\RG{29.743}{-1.684}{7.9701}{mi}{86}\RG{57.524}{-1.632}{7.9701}{rm}{49}\RG{61.759}{-1.632}{7.9701}{mi}{59}\RG{64.110}{-1.632}{7.9701}{mi}{116}\RG{150.955}{-1.632}{7.9701}{mi}{105}\RG{190.427}{-1.632}{7.9701}{mi}{106}\RG{86.573}{-13.083}{7.9701}{mi}{105}\RG{89.456}{-13.083}{7.9701}{rm}{61}\RG{96.039}{-13.083}{7.9701}{rm}{49}\RG{102.213}{-13.083}{7.9701}{mi}{106}\RG{106.097}{-13.083}{7.9701}{rm}{61}\RG{112.680}{-13.083}{7.9701}{rm}{49}\RG{89.680}{-24.231}{7.9701}{mi}{109}\RG{106.996}{-24.231}{7.9701}{mi}{110}\RG{86.455}{-27.504}{10.9092}{ex}{89}\RG{102.595}{-27.504}{10.9092}{ex}{89}\RG{116.918}{-29.031}{10.9092}{ex}{2}\RG{222.350}{-29.031}{10.9092}{ex}{3}\RG{0.000}{-37.868}{10.9091}{rm}{100}\RG{6.055}{-37.868}{10.9091}{rm}{101}\RG{10.898}{-37.868}{10.9091}{rm}{116}\RG{15.131}{-37.868}{10.9091}{rm}{40}\RG{19.394}{-37.868}{10.9091}{mi}{122}\RG{24.936}{-37.868}{10.9091}{mi}{73}\RG{39.478}{-37.868}{10.9091}{sy}{0}\RG{50.387}{-37.868}{10.9091}{mi}{39}\RG{67.667}{-37.868}{10.9091}{rm}{41}\RG{74.921}{-37.868}{10.9091}{rm}{61}\RG{121.463}{-37.868}{10.9091}{mi}{122}\RG{129.440}{-37.868}{10.9091}{sy}{0}\RG{140.349}{-37.868}{10.9091}{mi}{21}\RG{152.518}{-37.868}{10.9091}{rm}{43}\RG{163.428}{-37.868}{10.9091}{mi}{22}\RG{176.807}{-37.868}{10.9091}{sy}{0}\RG{187.717}{-37.868}{10.9091}{rm}{40}\RG{191.960}{-37.868}{10.9091}{mi}{105}\RG{198.142}{-37.868}{10.9091}{sy}{0}\RG{209.051}{-37.868}{10.9091}{mi}{106}\RG{214.168}{-37.868}{10.9091}{rm}{41}\RG{218.411}{-37.868}{10.9091}{mi}{116}\RG{226.895}{-37.868}{10.9091}{mi}{58}\RG{57.524}{-39.504}{7.9701}{rm}{50}\RG{61.758}{-39.504}{7.9701}{mi}{59}\RG{64.110}{-39.504}{7.9701}{mi}{116}\RG{146.713}{-39.504}{7.9701}{mi}{105}\RG{170.001}{-39.504}{7.9701}{mi}{106}\RG{29.743}{-39.556}{7.9701}{mi}{86}\RG{86.573}{-50.955}{7.9701}{mi}{105}\RG{89.456}{-50.955}{7.9701}{rm}{61}\RG{96.039}{-50.955}{7.9701}{rm}{49}\RG{102.213}{-50.955}{7.9701}{mi}{106}\RG{106.097}{-50.955}{7.9701}{rm}{61}\RG{112.680}{-50.955}{7.9701}{rm}{49}%
}
\]
在固定的标准基\allowbreak{}\hskip.1em \RB{14.722}{7.920}{3.760}{%
\RG{0.000}{0.002}{10.9091}{mi}{69}\RG{8.053}{-1.634}{7.9701}{mi}{105}\RG{10.930}{-1.634}{7.9701}{mi}{106}%
}\hskip.1em\allowbreak{} 下，这两个算子的矩阵元素均连续依赖于\allowbreak{}\hskip.1em \RB{4.100}{7.331}{0.618}{%
\RG{0.000}{0.002}{10.9091}{mi}{116}%
}\hskip.1em\allowbreak{} ，而行列式是矩阵元素的多项式。令\allowbreak{}\hskip.1em \RB{32.968}{9.270}{0.743}{%
\RG{0.000}{-0.003}{10.9091}{mi}{116}\RG{6.970}{-0.003}{10.9091}{sy}{33}\RG{20.909}{-0.003}{10.9091}{rm}{48}\RG{26.363}{3.956}{7.9701}{rm}{43}%
}\hskip.1em\allowbreak{} ，得到

\[
\RB{175.507}{55.519}{15.217}{%
\RG{84.270}{51.504}{7.9701}{mi}{109}\RG{101.586}{51.504}{7.9701}{mi}{110}\RG{81.045}{48.231}{10.9092}{ex}{89}\RG{97.185}{48.231}{10.9092}{ex}{89}\RG{0.000}{37.868}{10.9091}{rm}{100}\RG{6.055}{37.868}{10.9091}{rm}{101}\RG{10.898}{37.868}{10.9091}{rm}{116}\RG{15.131}{37.868}{10.9091}{rm}{40}\RG{19.394}{37.868}{10.9091}{mi}{122}\RG{24.947}{37.868}{10.9091}{mi}{73}\RG{39.479}{37.868}{10.9091}{sy}{0}\RG{50.388}{37.868}{10.9091}{mi}{39}\RG{62.256}{37.868}{10.9091}{rm}{41}\RG{69.511}{37.868}{10.9091}{rm}{61}\RG{111.507}{37.868}{10.9091}{rm}{40}\RG{115.750}{37.868}{10.9091}{mi}{122}\RG{123.727}{37.868}{10.9091}{sy}{0}\RG{134.636}{37.868}{10.9091}{mi}{21}\RG{144.381}{37.868}{10.9091}{mi}{22}\RG{155.337}{37.868}{10.9091}{rm}{41}\RG{159.579}{37.868}{10.9091}{mi}{59}\RG{57.524}{36.231}{7.9701}{rm}{49}\RG{141.000}{36.231}{7.9701}{mi}{105}\RG{150.954}{36.231}{7.9701}{mi}{106}\RG{29.743}{36.179}{7.9701}{mi}{86}\RG{81.163}{24.780}{7.9701}{mi}{105}\RG{84.046}{24.780}{7.9701}{rm}{61}\RG{90.629}{24.780}{7.9701}{rm}{49}\RG{96.803}{24.780}{7.9701}{mi}{106}\RG{100.687}{24.780}{7.9701}{rm}{61}\RG{107.270}{24.780}{7.9701}{rm}{49}\RG{84.270}{13.632}{7.9701}{mi}{109}\RG{101.586}{13.632}{7.9701}{mi}{110}\RG{81.045}{10.360}{10.9092}{ex}{89}\RG{97.185}{10.360}{10.9092}{ex}{89}\RG{0.000}{-0.004}{10.9091}{rm}{100}\RG{6.055}{-0.004}{10.9091}{rm}{101}\RG{10.898}{-0.004}{10.9091}{rm}{116}\RG{15.131}{-0.004}{10.9091}{rm}{40}\RG{19.394}{-0.004}{10.9091}{mi}{122}\RG{24.947}{-0.004}{10.9091}{mi}{73}\RG{39.479}{-0.004}{10.9091}{sy}{0}\RG{50.388}{-0.004}{10.9091}{mi}{39}\RG{62.256}{-0.004}{10.9091}{rm}{41}\RG{69.511}{-0.004}{10.9091}{rm}{61}\RG{111.507}{-0.004}{10.9091}{rm}{40}\RG{115.750}{-0.004}{10.9091}{mi}{122}\RG{123.727}{-0.004}{10.9091}{sy}{0}\RG{134.636}{-0.004}{10.9091}{mi}{21}\RG{146.805}{-0.004}{10.9091}{rm}{43}\RG{157.714}{-0.004}{10.9091}{mi}{22}\RG{168.670}{-0.004}{10.9091}{rm}{41}\RG{172.912}{-0.004}{10.9091}{mi}{58}\RG{57.524}{-1.640}{7.9701}{rm}{50}\RG{141.000}{-1.640}{7.9701}{mi}{105}\RG{164.288}{-1.640}{7.9701}{mi}{106}\RG{29.743}{-1.693}{7.9701}{mi}{86}\RG{81.163}{-13.091}{7.9701}{mi}{105}\RG{84.046}{-13.091}{7.9701}{rm}{61}\RG{90.629}{-13.091}{7.9701}{rm}{49}\RG{96.803}{-13.091}{7.9701}{mi}{106}\RG{100.687}{-13.091}{7.9701}{rm}{61}\RG{107.270}{-13.091}{7.9701}{rm}{49}%
}
\]
故两问的特征值分别为\allowbreak{}\hskip.1em \RB{20.111}{8.071}{3.762}{%
\RG{0.000}{-0.000}{10.9091}{mi}{21}\RG{6.364}{-1.636}{7.9701}{mi}{105}\RG{9.745}{-0.000}{10.9091}{mi}{22}\RG{16.319}{-1.636}{7.9701}{mi}{106}%
}\hskip.1em\allowbreak{} 与\allowbreak{}\hskip.1em \RB{148.575}{8.682}{3.762}{%
\RG{0.000}{-0.000}{10.9091}{mi}{21}\RG{6.364}{-1.636}{7.9701}{mi}{105}\RG{12.187}{-0.000}{10.9091}{sy}{0}\RG{23.113}{-0.000}{10.9091}{mi}{22}\RG{29.686}{-1.636}{7.9701}{mi}{106}\RG{37.728}{-0.000}{10.9091}{rm}{40}\RG{41.961}{-0.000}{10.9091}{rm}{49}\RG{50.499}{-0.000}{10.9091}{sy}{20}\RG{62.059}{-0.000}{10.9091}{mi}{105}\RG{68.891}{-0.000}{10.9091}{sy}{20}\RG{80.450}{-0.000}{10.9091}{mi}{109}\RG{90.028}{-0.000}{10.9091}{mi}{59}\RG{98.536}{-0.000}{10.9091}{rm}{49}\RG{107.065}{-0.000}{10.9091}{sy}{20}\RG{118.624}{-0.000}{10.9091}{mi}{106}\RG{126.815}{-0.000}{10.9091}{sy}{20}\RG{138.374}{-0.000}{10.9091}{mi}{110}\RG{144.922}{-0.000}{10.9091}{rm}{41}%
}\hskip.1em\allowbreak{} ，重复出现的次数就是相应的代数重数。\allowbreak{}\hskip.1em \RB{8.366}{8.000}{0.500}{%
\RG{0.000}{0.005}{10.9091}{am}{3}%
}\hskip.1em\allowbreak{}

方法小结：摄动后用秩一矩阵构造完整特征向量基，再对特征多项式取极限；不必要求各个乘积或差彼此不同。

题目来源：谢启鸿《复旦大学高等代数习题集（2023版）》，第115–116页。

% Original page 4.
\section*{2022S05}
证明：下列\allowbreak{}\hskip.1em \RB{6.729}{5.319}{0.623}{%
\RG{0.000}{-0.003}{10.9091}{mi}{110}%
}\hskip.1em\allowbreak{} 阶三对角矩阵可对角化：

\[
\RB{187.559}{26.413}{59.658}{%
\RG{22.727}{25.487}{10.9092}{ex}{48}\RG{173.601}{25.487}{10.9092}{ex}{49}\RG{42.213}{15.174}{10.9091}{rm}{48}\RG{77.505}{15.174}{10.9091}{rm}{49}\RG{22.727}{6.287}{10.9092}{ex}{66}\RG{173.601}{6.287}{10.9092}{ex}{67}\RG{32.273}{-0.001}{10.9091}{mi}{110}\RG{41.245}{-0.001}{10.9091}{sy}{0}\RG{52.154}{-0.001}{10.9091}{rm}{49}\RG{77.507}{-0.001}{10.9091}{rm}{48}\RG{106.504}{-0.001}{10.9091}{rm}{50}\RG{22.727}{-0.259}{10.9092}{ex}{66}\RG{173.601}{-0.259}{10.9092}{ex}{67}\RG{22.727}{-6.804}{10.9092}{ex}{66}\RG{173.601}{-6.804}{10.9092}{ex}{67}\RT{126.174}{-12.684}{10.9091}{.}\RG{22.727}{-13.350}{10.9092}{ex}{66}\RG{173.601}{-13.350}{10.9092}{ex}{67}\RT{130.419}{-15.673}{10.9091}{.}\RT{134.663}{-18.661}{10.9091}{.}\RG{0.000}{-19.350}{10.9091}{mi}{65}\RG{11.212}{-19.350}{10.9091}{rm}{61}\RG{184.964}{-19.350}{10.9091}{mi}{58}\RG{67.572}{-19.658}{10.9091}{mi}{110}\RG{76.544}{-19.658}{10.9091}{sy}{0}\RG{87.453}{-19.658}{10.9091}{rm}{50}\RG{106.500}{-19.658}{10.9091}{rm}{48}\RG{22.727}{-19.895}{10.9092}{ex}{66}\RG{173.601}{-19.895}{10.9092}{ex}{67}\RG{22.727}{-26.441}{10.9092}{ex}{66}\RG{173.601}{-26.441}{10.9092}{ex}{67}\RT{103.477}{-32.340}{10.9091}{.}\RT{126.174}{-32.340}{10.9091}{.}\RG{22.727}{-32.986}{10.9092}{ex}{66}\RG{173.601}{-32.986}{10.9092}{ex}{67}\RT{107.722}{-35.329}{10.9091}{.}\RT{130.419}{-35.329}{10.9091}{.}\RT{111.966}{-38.318}{10.9091}{.}\RT{134.663}{-38.318}{10.9091}{.}\RG{148.265}{-39.314}{10.9091}{mi}{110}\RG{157.237}{-39.314}{10.9091}{sy}{0}\RG{168.146}{-39.314}{10.9091}{rm}{49}\RG{22.727}{-39.968}{10.9092}{ex}{64}\RG{173.601}{-39.968}{10.9092}{ex}{65}\RG{129.208}{-54.489}{10.9091}{rm}{49}\RG{158.205}{-54.489}{10.9091}{rm}{48}%
}
\]
证明：当\allowbreak{}\hskip.1em \RB{26.147}{7.769}{0.617}{%
\RG{0.000}{0.003}{10.9091}{mi}{110}\RG{9.578}{0.003}{10.9091}{rm}{61}\RG{21.076}{0.003}{10.9091}{rm}{49}%
}\hskip.1em\allowbreak{} 时结论成立。以下设\allowbreak{}\hskip.1em \RB{26.492}{7.769}{1.992}{%
\RG{0.000}{0.003}{10.9091}{mi}{110}\RG{9.578}{0.003}{10.9091}{sy}{21}\RG{21.094}{0.003}{10.9091}{rm}{50}%
}\hskip.1em\allowbreak{} 。令

\[
\RB{267.903}{20.616}{13.112}{%
\RG{23.604}{19.679}{10.9092}{ex}{115}\RG{34.513}{15.383}{10.9092}{ex}{18}\RG{67.881}{15.383}{10.9092}{ex}{19}\RG{42.544}{7.381}{10.9091}{mi}{110}\RG{51.516}{7.381}{10.9091}{sy}{0}\RG{62.426}{7.381}{10.9091}{rm}{49}\RG{0.000}{0.001}{10.9091}{mi}{100}\RG{12.089}{0.001}{10.9091}{rm}{61}\RG{86.820}{0.001}{10.9091}{rm}{40}\RG{91.053}{0.001}{10.9091}{rm}{49}\RG{99.547}{0.001}{10.9091}{sy}{20}\RG{111.062}{0.001}{10.9091}{mi}{105}\RG{117.851}{0.001}{10.9091}{sy}{20}\RG{129.366}{0.001}{10.9091}{mi}{110}\RG{135.914}{0.001}{10.9091}{rm}{41}\RG{140.157}{0.001}{10.9091}{mi}{59}\RG{166.808}{0.001}{10.9091}{mi}{68}\RG{179.189}{0.001}{10.9091}{rm}{61}\RG{190.687}{0.001}{10.9091}{rm}{100}\RG{196.741}{0.001}{10.9091}{rm}{105}\RG{199.763}{0.001}{10.9091}{rm}{97}\RG{205.218}{0.001}{10.9091}{rm}{103}\RG{210.825}{0.001}{10.9091}{rm}{40}\RG{215.098}{0.001}{10.9091}{mi}{100}\RG{225.509}{0.001}{10.9091}{mi}{59}\RG{230.342}{0.001}{10.9091}{mi}{58}\RG{235.185}{0.001}{10.9091}{mi}{58}\RG{240.018}{0.001}{10.9091}{mi}{58}\RG{244.862}{0.001}{10.9091}{mi}{59}\RG{249.705}{0.001}{10.9091}{mi}{100}\RG{261.065}{0.001}{10.9091}{rm}{41}\RG{265.308}{0.001}{10.9091}{mi}{58}\RG{5.678}{-1.636}{7.9701}{mi}{105}\RG{220.776}{-1.636}{7.9701}{rm}{49}\RG{255.429}{-1.636}{7.9701}{mi}{110}\RG{43.939}{-7.482}{10.9091}{mi}{105}\RG{50.121}{-7.483}{10.9091}{sy}{0}\RG{61.031}{-7.482}{10.9091}{rm}{49}\RR{34.513}{19.897}{41.397}{0.436}%
}
\]
各个\allowbreak{}\hskip.1em \RB{8.705}{8.074}{2.213}{%
\RG{0.000}{0.003}{10.9091}{mi}{100}\RG{5.678}{-1.633}{7.9701}{mi}{105}%
}\hskip.1em\allowbreak{} 均为正数，且

\[
\RB{150.311}{18.156}{9.705}{%
\RG{35.621}{17.219}{10.9092}{ex}{114}\RG{0.000}{7.375}{10.9091}{mi}{100}\RG{47.725}{7.375}{10.9091}{mi}{110}\RG{56.697}{7.375}{10.9091}{sy}{0}\RG{67.608}{7.375}{10.9091}{mi}{105}\RG{5.678}{5.739}{7.9701}{mi}{105}\RG{8.561}{5.739}{7.9701}{rm}{43}\RG{15.144}{5.739}{7.9701}{rm}{49}\RG{24.106}{-0.005}{10.9091}{rm}{61}\RG{94.380}{-0.005}{10.9091}{rm}{40}\RG{98.613}{-0.005}{10.9091}{rm}{49}\RG{107.107}{-0.005}{10.9091}{sy}{20}\RG{118.622}{-0.005}{10.9091}{mi}{105}\RG{125.397}{-0.005}{10.9091}{mi}{60}\RG{136.906}{-0.005}{10.9091}{mi}{110}\RG{143.474}{-0.005}{10.9091}{rm}{41}\RG{147.716}{-0.005}{10.9091}{mi}{58}\RG{5.410}{-7.488}{10.9091}{mi}{100}\RG{57.666}{-7.488}{10.9091}{mi}{105}\RG{11.088}{-9.125}{7.9701}{mi}{105}\RR{-0.000}{2.722}{19.880}{0.436}\RR{46.530}{17.437}{26.030}{0.436}\RR{47.725}{2.722}{23.640}{0.436}%
}
\]
记\allowbreak{}\hskip.1em \RB{62.606}{9.757}{0.502}{%
\RG{0.000}{-0.002}{10.9091}{mi}{72}\RG{12.985}{-0.002}{10.9091}{rm}{61}\RG{24.500}{-0.002}{10.9091}{mi}{68}\RG{33.835}{3.957}{7.9701}{sy}{0}\RG{40.421}{3.957}{7.9701}{rm}{49}\RG{45.153}{-0.002}{10.9091}{mi}{65}\RG{53.335}{-0.002}{10.9091}{mi}{68}%
}\hskip.1em\allowbreak{} 。它仍是对角元为零的三对角阵，并且

\[
\RB{158.597}{43.513}{10.447}{%
\RG{80.844}{37.922}{10.9092}{ex}{112}\RG{45.222}{35.442}{10.9091}{mi}{100}\RG{50.900}{33.806}{7.9701}{mi}{105}\RG{53.783}{33.806}{7.9701}{rm}{43}\RG{60.367}{33.806}{7.9701}{rm}{49}\RG{0.000}{28.062}{10.9091}{mi}{104}\RG{28.753}{28.062}{10.9091}{rm}{61}\RG{40.268}{28.062}{10.9091}{mi}{105}\RG{69.329}{28.062}{10.9091}{rm}{61}\RG{91.753}{28.062}{10.9091}{mi}{105}\RG{95.511}{28.062}{10.9091}{rm}{40}\RG{99.754}{28.062}{10.9091}{mi}{110}\RG{108.726}{28.062}{10.9091}{sy}{0}\RG{119.635}{28.062}{10.9091}{mi}{105}\RG{123.394}{28.062}{10.9091}{rm}{41}\RG{127.636}{28.062}{10.9091}{mi}{59}\RG{6.286}{26.426}{7.9701}{mi}{105}\RG{9.163}{26.426}{7.9701}{mi}{59}\RG{11.514}{26.426}{7.9701}{mi}{105}\RG{14.404}{26.426}{7.9701}{rm}{43}\RG{20.987}{26.426}{7.9701}{rm}{49}\RG{50.633}{20.579}{10.9091}{mi}{100}\RG{56.311}{18.942}{7.9701}{mi}{105}\RG{109.210}{9.861}{10.9092}{ex}{112}\RG{79.000}{7.382}{10.9091}{mi}{100}\RG{84.678}{5.745}{7.9701}{mi}{105}\RG{0.000}{0.002}{10.9091}{mi}{104}\RG{28.753}{0.002}{10.9091}{rm}{61}\RG{40.252}{0.002}{10.9091}{rm}{40}\RG{44.511}{0.002}{10.9091}{mi}{110}\RG{53.483}{0.002}{10.9091}{sy}{0}\RG{64.393}{0.002}{10.9091}{mi}{105}\RG{68.151}{0.002}{10.9091}{rm}{41}\RG{97.695}{0.002}{10.9091}{rm}{61}\RG{120.119}{0.002}{10.9091}{mi}{105}\RG{123.878}{0.002}{10.9091}{rm}{40}\RG{128.120}{0.002}{10.9091}{mi}{110}\RG{137.092}{0.002}{10.9091}{sy}{0}\RG{148.001}{0.002}{10.9091}{mi}{105}\RG{151.760}{0.002}{10.9091}{rm}{41}\RG{156.002}{0.002}{10.9091}{mi}{58}\RG{6.286}{-1.635}{7.9701}{mi}{105}\RG{9.169}{-1.635}{7.9701}{rm}{43}\RG{15.752}{-1.635}{7.9701}{rm}{49}\RG{19.989}{-1.635}{7.9701}{mi}{59}\RG{22.340}{-1.635}{7.9701}{mi}{105}\RG{73.589}{-7.482}{10.9091}{mi}{100}\RG{79.267}{-9.118}{7.9701}{mi}{105}\RG{82.151}{-9.118}{7.9701}{rm}{43}\RG{88.734}{-9.118}{7.9701}{rm}{49}\RR{45.222}{30.789}{19.880}{0.436}\RR{91.753}{38.140}{35.883}{0.436}\RR{73.589}{2.729}{19.880}{0.436}\RR{120.119}{10.080}{35.883}{0.436}%
}
\]
故\allowbreak{}\hskip.1em \RB{10.111}{7.954}{0.500}{%
\RG{0.000}{0.003}{10.9091}{mi}{72}%
}\hskip.1em\allowbreak{} 为实对称阵，存在实正交矩阵\allowbreak{}\hskip.1em \RB{8.573}{8.194}{2.613}{%
\RG{0.000}{0.003}{10.9091}{mi}{81}%
}\hskip.1em\allowbreak{} 使\allowbreak{}\hskip.1em \RB{55.725}{9.938}{2.613}{%
\RG{0.000}{0.003}{10.9091}{mi}{81}\RG{8.624}{3.962}{7.9701}{rm}{84}\RG{14.906}{0.003}{10.9091}{mi}{72}\RG{24.855}{0.003}{10.9091}{mi}{81}\RG{36.515}{0.003}{10.9091}{rm}{61}\RG{48.014}{0.003}{10.9091}{rm}{3}%
}\hskip.1em\allowbreak{} 为实对角阵。于是

\[
\RB{173.472}{10.479}{3.228}{%
\RG{26.444}{4.503}{7.9701}{sy}{0}\RG{33.030}{4.503}{7.9701}{rm}{49}\RG{95.559}{4.503}{7.9701}{rm}{84}\RG{111.176}{4.503}{7.9701}{sy}{0}\RG{117.762}{4.503}{7.9701}{rm}{49}\RG{0.000}{-0.001}{10.9091}{rm}{40}\RG{4.242}{-0.001}{10.9091}{mi}{68}\RG{13.558}{-0.001}{10.9091}{mi}{81}\RG{22.201}{-0.001}{10.9091}{rm}{41}\RG{37.762}{-0.001}{10.9091}{mi}{65}\RG{45.944}{-0.001}{10.9091}{rm}{40}\RG{50.187}{-0.001}{10.9091}{mi}{68}\RG{59.504}{-0.001}{10.9091}{mi}{81}\RG{68.147}{-0.001}{10.9091}{rm}{41}\RG{75.401}{-0.001}{10.9091}{rm}{61}\RG{86.935}{-0.001}{10.9091}{mi}{81}\RG{101.841}{-0.001}{10.9091}{mi}{68}\RG{122.495}{-0.001}{10.9091}{mi}{65}\RG{130.677}{-0.001}{10.9091}{mi}{68}\RG{139.993}{-0.001}{10.9091}{mi}{81}\RG{151.666}{-0.001}{10.9091}{rm}{61}\RG{163.165}{-0.001}{10.9091}{rm}{3}\RG{170.757}{-0.001}{10.9091}{mi}{59}%
}
\]
即\allowbreak{}\hskip.1em \RB{8.366}{8.316}{0.500}{%
\RG{0.000}{0.005}{10.9091}{mi}{65}%
}\hskip.1em\allowbreak{} 在实数域上可对角化。

这里所用的实对称阵对角化结论可简证如下：对称阵\allowbreak{}\hskip.1em \RB{10.111}{7.952}{0.500}{%
\RG{0.000}{0.001}{10.9091}{mi}{72}%
}\hskip.1em\allowbreak{} 的实二次型\allowbreak{}\hskip.1em \RB{28.715}{9.936}{0.619}{%
\RG{0.000}{0.001}{10.9091}{mi}{120}\RG{6.235}{3.960}{7.9701}{rm}{84}\RG{12.517}{0.001}{10.9091}{mi}{72}\RG{22.466}{0.001}{10.9091}{mi}{120}%
}\hskip.1em\allowbreak{} 在单位球面上取得最大值。若最大值点为\allowbreak{}\hskip.1em \RB{5.606}{5.318}{0.624}{%
\RG{0.000}{-0.004}{10.9091}{mi}{118}%
}\hskip.1em\allowbreak{} ，对任意\allowbreak{}\hskip.1em \RB{28.677}{7.783}{0.624}{%
\RG{0.000}{-0.004}{10.9091}{mi}{119}\RG{11.345}{-0.004}{10.9091}{sy}{63}\RG{23.071}{-0.004}{10.9091}{mi}{118}%
}\hskip.1em\allowbreak{} ，考察\allowbreak{}\hskip.1em \RB{86.347}{8.678}{3.231}{%
\RG{0.000}{-0.004}{10.9091}{rm}{40}\RG{4.243}{-0.004}{10.9091}{mi}{118}\RG{12.431}{-0.004}{10.9091}{rm}{43}\RG{23.424}{-0.004}{10.9091}{mi}{116}\RG{27.362}{-0.004}{10.9091}{mi}{119}\RG{35.467}{-0.004}{10.9091}{rm}{41}\RG{39.710}{-0.004}{10.9091}{mi}{61}\RG{45.164}{-0.004}{10.9091}{sy}{107}\RG{50.619}{-0.004}{10.9091}{mi}{118}\RG{58.807}{-0.004}{10.9091}{rm}{43}\RG{69.801}{-0.004}{10.9091}{mi}{116}\RG{73.739}{-0.004}{10.9091}{mi}{119}\RG{81.844}{-0.004}{10.9091}{sy}{107}%
}\hskip.1em\allowbreak{} 处的值并在\allowbreak{}\hskip.1em \RB{24.416}{7.761}{0.744}{%
\RG{0.000}{-0.004}{10.9091}{mi}{116}\RG{7.181}{-0.004}{10.9091}{rm}{61}\RG{18.898}{-0.004}{10.9091}{rm}{48}%
}\hskip.1em\allowbreak{} 求导，得到\allowbreak{}\hskip.1em \RB{50.066}{9.936}{0.739}{%
\RG{0.000}{0.001}{10.9091}{mi}{119}\RG{8.103}{3.960}{7.9701}{rm}{84}\RG{14.385}{0.001}{10.9091}{mi}{72}\RG{24.334}{0.001}{10.9091}{mi}{118}\RG{33.049}{0.001}{10.9091}{rm}{61}\RG{44.548}{0.001}{10.9091}{rm}{48}%
}\hskip.1em\allowbreak{} ，故\allowbreak{}\hskip.1em \RB{15.555}{7.952}{0.619}{%
\RG{0.000}{0.001}{10.9091}{mi}{72}\RG{9.949}{0.001}{10.9091}{mi}{118}%
}\hskip.1em\allowbreak{} 平行于\allowbreak{}\hskip.1em \RB{5.606}{5.323}{0.619}{%
\RG{0.000}{0.001}{10.9091}{mi}{118}%
}\hskip.1em\allowbreak{} 。又由对称性，\allowbreak{}\hskip.1em \RB{12.292}{9.991}{0.619}{%
\RG{0.000}{0.001}{10.9091}{mi}{118}\RG{5.679}{3.960}{7.9701}{sy}{63}%
}\hskip.1em\allowbreak{} 在\allowbreak{}\hskip.1em \RB{10.111}{7.952}{0.500}{%
\RG{0.000}{0.001}{10.9091}{mi}{72}%
}\hskip.1em\allowbreak{} 下不变；对维数归纳，即得一组标准正交特征向量基。\allowbreak{}\hskip.1em \RB{8.366}{7.991}{0.504}{%
\RG{0.000}{-0.004}{10.9091}{am}{3}%
}\hskip.1em\allowbreak{}

方法小结：对角相似变换把对应的上下副对角元配成相等的正数，从而将问题归结为实对称阵的对角化。

题目来源：谢启鸿《复旦大学高等代数习题集（2023版）》，第116页。

% Original page 5.
\section*{2022S06}
设\allowbreak{}\hskip.1em \RB{56.468}{8.679}{3.230}{%
\RG{0.000}{-0.003}{10.9091}{mi}{65}\RG{11.212}{-0.003}{10.9091}{sy}{50}\RG{21.515}{-0.003}{10.9091}{mi}{77}\RG{32.099}{-1.639}{7.9701}{mi}{109}\RG{40.087}{-0.003}{10.9091}{rm}{40}\RG{44.330}{-0.003}{10.9091}{bb}{75}\RG{52.815}{-0.003}{10.9091}{rm}{41}%
}\hskip.1em\allowbreak{} 、\allowbreak{}\hskip.1em \RB{54.755}{8.679}{3.230}{%
\RG{0.000}{-0.003}{10.9091}{mi}{66}\RG{11.852}{-0.003}{10.9091}{sy}{50}\RG{22.155}{-0.003}{10.9091}{mi}{77}\RG{32.738}{-1.639}{7.9701}{mi}{110}\RG{38.375}{-0.003}{10.9091}{rm}{40}\RG{42.617}{-0.003}{10.9091}{bb}{75}\RG{51.102}{-0.003}{10.9091}{rm}{41}%
}\hskip.1em\allowbreak{} ，其中\allowbreak{}\hskip.1em \RB{8.889}{7.970}{0.503}{%
\RG{0.000}{-0.003}{10.9091}{bb}{75}%
}\hskip.1em\allowbreak{} 为数域。若\allowbreak{}\hskip.1em \RB{21.762}{8.308}{2.608}{%
\RG{0.000}{-0.003}{10.9091}{mi}{65}\RG{8.182}{-0.003}{10.9091}{mi}{59}\RG{13.015}{-0.003}{10.9091}{mi}{66}%
}\hskip.1em\allowbreak{} 在复数域中有公共特征值，证明：存在非零矩阵\allowbreak{}\hskip.1em \RB{68.588}{8.684}{3.225}{%
\RG{0.000}{0.002}{10.9091}{mi}{67}\RG{11.607}{0.002}{10.9091}{sy}{50}\RG{21.910}{0.002}{10.9091}{mi}{77}\RG{32.494}{-1.634}{7.9701}{mi}{109}\RG{39.984}{-1.634}{7.9701}{sy}{2}\RG{46.571}{-1.634}{7.9701}{mi}{110}\RG{52.207}{0.002}{10.9091}{rm}{40}\RG{56.450}{0.002}{10.9091}{bb}{75}\RG{64.936}{0.002}{10.9091}{rm}{41}%
}\hskip.1em\allowbreak{} ，使得\allowbreak{}\hskip.1em \RB{48.615}{8.313}{0.738}{%
\RG{0.000}{0.002}{10.9091}{mi}{65}\RG{8.182}{0.002}{10.9091}{mi}{67}\RG{19.789}{0.002}{10.9091}{rm}{61}\RG{31.304}{0.002}{10.9091}{mi}{67}\RG{39.868}{0.002}{10.9091}{mi}{66}%
}\hskip.1em\allowbreak{} 。

证明：先在复数域中构造一个非零解：设\allowbreak{}\hskip.1em \RB{6.478}{8.070}{0.643}{%
\RG{0.000}{-0.001}{10.9091}{mi}{21}%
}\hskip.1em\allowbreak{} 为公共特征值。由于转置不改变特征多项式，可取非零列向量\allowbreak{}\hskip.1em \RB{35.186}{8.183}{0.933}{%
\RG{0.000}{0.003}{10.9091}{mi}{117}\RG{9.275}{0.003}{10.9091}{sy}{50}\RG{19.578}{0.003}{10.9091}{bb}{67}\RG{27.457}{3.962}{7.9701}{mi}{109}%
}\hskip.1em\allowbreak{} 、\allowbreak{}\hskip.1em \RB{32.269}{8.183}{0.933}{%
\RG{0.000}{0.003}{10.9091}{mi}{118}\RG{8.709}{0.003}{10.9091}{sy}{50}\RG{19.012}{0.003}{10.9091}{bb}{67}\RG{26.891}{3.962}{7.9701}{mi}{110}%
}\hskip.1em\allowbreak{} ，使

\[
\RB{118.193}{10.479}{2.607}{%
\RG{77.071}{4.503}{7.9701}{rm}{84}\RG{0.000}{-0.002}{10.9091}{mi}{65}\RG{8.182}{-0.002}{10.9091}{mi}{117}\RG{17.457}{-0.002}{10.9091}{rm}{61}\RG{28.972}{-0.002}{10.9091}{mi}{21}\RG{35.332}{-0.002}{10.9091}{mi}{117}\RG{41.572}{-0.002}{10.9091}{mi}{59}\RG{68.223}{-0.002}{10.9091}{mi}{66}\RG{83.353}{-0.002}{10.9091}{mi}{118}\RG{92.062}{-0.002}{10.9091}{rm}{61}\RG{103.577}{-0.002}{10.9091}{mi}{21}\RG{109.937}{-0.002}{10.9091}{mi}{118}\RG{115.599}{-0.002}{10.9091}{mi}{58}%
}
\]
令\allowbreak{}\hskip.1em \RB{44.975}{9.933}{2.306}{%
\RG{0.000}{-0.002}{10.9091}{mi}{67}\RG{7.797}{-1.639}{7.9701}{rm}{48}\RG{15.559}{-0.002}{10.9091}{rm}{61}\RG{27.075}{-0.002}{10.9091}{mi}{117}\RG{33.315}{-0.002}{10.9091}{mi}{118}\RG{38.999}{3.957}{7.9701}{rm}{84}%
}\hskip.1em\allowbreak{} 。选择非零分量\allowbreak{}\hskip.1em \RB{23.555}{5.320}{3.765}{%
\RG{0.000}{-0.002}{10.9091}{mi}{117}\RG{6.245}{-1.639}{7.9701}{mi}{105}\RG{9.626}{-0.002}{10.9091}{mi}{59}\RG{14.459}{-0.002}{10.9091}{mi}{118}\RG{19.763}{-1.639}{7.9701}{mi}{106}%
}\hskip.1em\allowbreak{} ，可知\allowbreak{}\hskip.1em \RB{82.168}{8.680}{3.765}{%
\RG{0.000}{-0.002}{10.9091}{rm}{40}\RG{4.242}{-0.002}{10.9091}{mi}{67}\RG{12.039}{-1.639}{7.9701}{rm}{48}\RG{16.772}{-0.002}{10.9091}{rm}{41}\RG{21.014}{-1.639}{7.9701}{mi}{105}\RG{23.891}{-1.639}{7.9701}{mi}{106}\RG{31.310}{-0.002}{10.9091}{rm}{61}\RG{42.825}{-0.002}{10.9091}{mi}{117}\RG{49.070}{-1.639}{7.9701}{mi}{105}\RG{52.451}{-0.002}{10.9091}{mi}{118}\RG{57.739}{-1.639}{7.9701}{mi}{106}\RG{65.152}{-0.002}{10.9091}{sy}{54}\RG{65.152}{-0.002}{10.9091}{rm}{61}\RG{76.650}{-0.002}{10.9091}{rm}{48}%
}\hskip.1em\allowbreak{} ，所以\allowbreak{}\hskip.1em \RB{32.577}{8.309}{2.858}{%
\RG{0.000}{-0.002}{10.9091}{mi}{67}\RG{7.797}{-1.639}{7.9701}{rm}{48}\RG{15.560}{-0.002}{10.9091}{sy}{54}\RG{15.560}{-0.002}{10.9091}{rm}{61}\RG{27.058}{-0.002}{10.9091}{rm}{48}%
}\hskip.1em\allowbreak{} 。而且

\[
\RB{204.071}{10.484}{3.223}{%
\RG{63.848}{4.508}{7.9701}{rm}{84}\RG{102.964}{4.508}{7.9701}{rm}{84}\RG{143.100}{4.508}{7.9701}{rm}{84}\RG{159.304}{4.508}{7.9701}{rm}{84}\RG{0.000}{0.004}{10.9091}{mi}{65}\RG{8.182}{0.004}{10.9091}{mi}{67}\RG{23.741}{0.004}{10.9091}{rm}{61}\RG{35.239}{0.004}{10.9091}{rm}{40}\RG{39.499}{0.004}{10.9091}{mi}{65}\RG{47.681}{0.004}{10.9091}{mi}{117}\RG{53.925}{0.004}{10.9091}{rm}{41}\RG{58.168}{0.004}{10.9091}{mi}{118}\RG{73.161}{0.004}{10.9091}{rm}{61}\RG{84.676}{0.004}{10.9091}{mi}{21}\RG{91.036}{0.004}{10.9091}{mi}{117}\RG{97.276}{0.004}{10.9091}{mi}{118}\RG{112.276}{0.004}{10.9091}{rm}{61}\RG{123.791}{0.004}{10.9091}{mi}{117}\RG{130.036}{0.004}{10.9091}{rm}{40}\RG{134.278}{0.004}{10.9091}{mi}{66}\RG{149.383}{0.004}{10.9091}{mi}{118}\RG{155.062}{0.004}{10.9091}{rm}{41}\RG{168.617}{0.004}{10.9091}{rm}{61}\RG{180.132}{0.004}{10.9091}{mi}{67}\RG{192.662}{0.004}{10.9091}{mi}{66}\RG{201.476}{0.004}{10.9091}{mi}{58}\RG{15.978}{-1.633}{7.9701}{rm}{48}\RG{187.929}{-1.633}{7.9701}{rm}{48}%
}
\]
注意这里使用普通转置，不使用共轭转置。

现在将解降回\allowbreak{}\hskip.1em \RB{8.889}{7.969}{0.504}{%
\RG{0.000}{-0.004}{10.9091}{bb}{75}%
}\hskip.1em\allowbreak{} 。设未知矩阵\allowbreak{}\hskip.1em \RB{43.004}{8.678}{3.766}{%
\RG{0.000}{-0.004}{10.9091}{mi}{67}\RG{11.607}{-0.004}{10.9091}{rm}{61}\RG{23.105}{-0.004}{10.9091}{rm}{40}\RG{27.365}{-0.004}{10.9091}{mi}{99}\RG{32.086}{-1.640}{7.9701}{mi}{105}\RG{34.963}{-1.640}{7.9701}{mi}{106}\RG{39.351}{-0.004}{10.9091}{rm}{41}%
}\hskip.1em\allowbreak{} ，则方程\allowbreak{}\hskip.1em \RB{48.616}{8.307}{0.744}{%
\RG{0.000}{-0.004}{10.9091}{mi}{65}\RG{8.182}{-0.004}{10.9091}{mi}{67}\RG{19.790}{-0.004}{10.9091}{rm}{61}\RG{31.305}{-0.004}{10.9091}{mi}{67}\RG{39.869}{-0.004}{10.9091}{mi}{66}%
}\hskip.1em\allowbreak{} 等价于

\[
\RB{259.174}{17.647}{13.953}{%
\RG{4.133}{13.632}{7.9701}{mi}{109}\RG{63.714}{13.632}{7.9701}{mi}{110}\RG{0.000}{10.360}{10.9092}{ex}{88}\RG{58.403}{10.360}{10.9092}{ex}{88}\RG{17.576}{-0.004}{10.9091}{mi}{97}\RG{31.345}{-0.004}{10.9091}{mi}{99}\RG{47.494}{-0.004}{10.9091}{sy}{0}\RG{75.980}{-0.004}{10.9091}{mi}{99}\RG{87.541}{-0.004}{10.9091}{mi}{98}\RG{103.095}{-0.004}{10.9091}{rm}{61}\RG{114.593}{-0.004}{10.9091}{rm}{48}\RG{120.064}{-0.004}{10.9091}{mi}{59}\RG{146.731}{-0.004}{10.9091}{rm}{49}\RG{155.216}{-0.004}{10.9091}{sy}{20}\RG{166.731}{-0.004}{10.9091}{mi}{105}\RG{173.520}{-0.004}{10.9091}{sy}{20}\RG{185.036}{-0.004}{10.9091}{mi}{109}\RG{194.614}{-0.004}{10.9091}{mi}{59}\RG{210.372}{-0.004}{10.9091}{rm}{49}\RG{218.856}{-0.004}{10.9091}{sy}{20}\RG{230.371}{-0.004}{10.9091}{mi}{106}\RG{238.519}{-0.004}{10.9091}{sy}{20}\RG{250.034}{-0.004}{10.9091}{mi}{110}\RG{256.580}{-0.004}{10.9091}{mi}{58}\RG{23.342}{-1.781}{7.9701}{mi}{105}\RG{26.219}{-1.781}{7.9701}{mi}{107}\RG{36.066}{-1.781}{7.9701}{mi}{107}\RG{40.681}{-1.781}{7.9701}{mi}{106}\RG{80.701}{-1.781}{7.9701}{mi}{105}\RG{83.578}{-1.781}{7.9701}{mi}{96}\RG{92.223}{-1.781}{7.9701}{mi}{96}\RG{95.682}{-1.781}{7.9701}{mi}{106}\RG{0.157}{-13.357}{7.9701}{mi}{107}\RG{4.779}{-13.357}{7.9701}{rm}{61}\RG{11.362}{-13.357}{7.9701}{rm}{49}\RG{59.142}{-13.357}{7.9701}{mi}{96}\RG{62.601}{-13.357}{7.9701}{rm}{61}\RG{69.185}{-13.357}{7.9701}{rm}{49}%
}
\]
将\allowbreak{}\hskip.1em \RB{16.307}{5.327}{0.615}{%
\RG{0.000}{0.005}{10.9091}{mi}{109}\RG{9.578}{0.005}{10.9091}{mi}{110}%
}\hskip.1em\allowbreak{} 个未知数按固定顺序排成列向量\allowbreak{}\hskip.1em \RB{5.191}{5.327}{0.615}{%
\RG{0.000}{0.005}{10.9091}{mi}{99}%
}\hskip.1em\allowbreak{} ，便得到齐次线性方程组

\[
\RB{131.542}{8.683}{3.226}{%
\RG{0.000}{0.001}{10.9091}{mi}{77}\RG{11.760}{0.001}{10.9091}{mi}{99}\RG{19.524}{0.001}{10.9091}{rm}{61}\RG{31.022}{0.001}{10.9091}{rm}{48}\RG{36.493}{0.001}{10.9091}{mi}{59}\RG{63.144}{0.001}{10.9091}{mi}{77}\RG{77.964}{0.001}{10.9091}{sy}{50}\RG{88.267}{0.001}{10.9091}{mi}{77}\RG{111.977}{0.001}{10.9091}{rm}{40}\RG{116.220}{0.001}{10.9091}{bb}{75}\RG{124.704}{0.001}{10.9091}{rm}{41}\RG{128.947}{0.001}{10.9091}{mi}{58}\RG{98.850}{-1.635}{7.9701}{mi}{109}\RG{106.334}{-1.635}{7.9701}{mi}{110}%
}
\]
上一步构造的\allowbreak{}\hskip.1em \RB{12.194}{8.188}{2.306}{%
\RG{0.000}{-0.003}{10.9091}{mi}{67}\RG{7.797}{-1.639}{7.9701}{rm}{48}%
}\hskip.1em\allowbreak{} 表明，这个方程组在\allowbreak{}\hskip.1em \RB{7.973}{8.177}{0.710}{%
\RG{0.000}{-0.003}{10.9091}{bb}{67}%
}\hskip.1em\allowbreak{} 上有非零解，故\allowbreak{}\hskip.1em \RB{48.789}{8.068}{0.743}{%
\RG{0.000}{-0.003}{10.9091}{rm}{100}\RG{6.055}{-0.003}{10.9091}{rm}{101}\RG{10.898}{-0.003}{10.9091}{rm}{116}\RG{16.969}{-0.003}{10.9091}{mi}{77}\RG{31.772}{-0.003}{10.9091}{rm}{61}\RG{43.271}{-0.003}{10.9091}{rm}{48}%
}\hskip.1em\allowbreak{} 。行列式由\allowbreak{}\hskip.1em \RB{11.889}{7.948}{0.503}{%
\RG{0.000}{-0.003}{10.9091}{mi}{77}%
}\hskip.1em\allowbreak{} 的元素通过加法和乘法计算，其在\allowbreak{}\hskip.1em \RB{8.889}{7.975}{0.500}{%
\RG{0.000}{0.002}{10.9091}{bb}{75}%
}\hskip.1em\allowbreak{} 与\allowbreak{}\hskip.1em \RB{7.973}{8.182}{0.705}{%
\RG{0.000}{0.002}{10.9091}{bb}{67}%
}\hskip.1em\allowbreak{} 中的值一致。因此\allowbreak{}\hskip.1em \RB{11.889}{7.953}{0.500}{%
\RG{0.000}{0.002}{10.9091}{mi}{77}%
}\hskip.1em\allowbreak{} 在\allowbreak{}\hskip.1em \RB{8.889}{7.975}{0.500}{%
\RG{0.000}{0.002}{10.9091}{bb}{75}%
}\hskip.1em\allowbreak{} 上也不可逆。

在\allowbreak{}\hskip.1em \RB{8.889}{7.971}{0.502}{%
\RG{0.000}{-0.002}{10.9091}{bb}{75}%
}\hskip.1em\allowbreak{} 上对\allowbreak{}\hskip.1em \RB{11.889}{7.949}{0.502}{%
\RG{0.000}{-0.002}{10.9091}{mi}{77}%
}\hskip.1em\allowbreak{} 作高斯消元，主元个数少于\allowbreak{}\hskip.1em \RB{16.307}{5.320}{0.622}{%
\RG{0.000}{-0.002}{10.9091}{mi}{109}\RG{9.578}{-0.002}{10.9091}{mi}{110}%
}\hskip.1em\allowbreak{} ，因而存在自由未知数；将其中一个自由未知数取为\allowbreak{}\hskip.1em \RB{5.071}{7.769}{0.500}{%
\RG{0.000}{0.003}{10.9091}{rm}{49}%
}\hskip.1em\allowbreak{} ，其余取为\allowbreak{}\hskip.1em \RB{5.518}{7.769}{0.737}{%
\RG{0.000}{0.003}{10.9091}{rm}{48}%
}\hskip.1em\allowbreak{} ，回代即可得到\allowbreak{}\hskip.1em \RB{21.347}{7.977}{0.500}{%
\RG{0.000}{0.003}{10.9091}{bb}{75}\RG{8.485}{3.962}{7.9701}{mi}{109}\RG{15.969}{3.962}{7.9701}{mi}{110}%
}\hskip.1em\allowbreak{} 中的非零解：将它还原为矩阵\allowbreak{}\hskip.1em \RB{8.791}{8.194}{0.737}{%
\RG{0.000}{0.003}{10.9091}{mi}{67}%
}\hskip.1em\allowbreak{} ，即得所求。\allowbreak{}\hskip.1em \RB{8.366}{7.998}{0.500}{%
\RG{0.000}{0.003}{10.9091}{am}{3}%
}\hskip.1em\allowbreak{}

方法小结：先在扩域中用左右特征向量构造解，再利用系数矩阵的奇异性回到原数域；不要求那些复特征向量本身属于\allowbreak{}\hskip.1em \RB{8.889}{7.970}{0.503}{%
\RG{0.000}{-0.003}{10.9091}{bb}{75}%
}\hskip.1em\allowbreak{} 。

题目来源：谢启鸿《复旦大学高等代数习题集（2023版）》，第116页。

% Original page 6.
\section*{2022S07}
设\allowbreak{}\hskip.1em \RB{54.115}{8.679}{3.230}{%
\RG{0.000}{-0.003}{10.9091}{mi}{65}\RG{11.212}{-0.003}{10.9091}{sy}{50}\RG{21.515}{-0.003}{10.9091}{mi}{77}\RG{32.098}{-1.639}{7.9701}{mi}{110}\RG{37.735}{-0.003}{10.9091}{rm}{40}\RG{41.977}{-0.003}{10.9091}{bb}{75}\RG{50.462}{-0.003}{10.9091}{rm}{41}%
}\hskip.1em\allowbreak{} ，其中\allowbreak{}\hskip.1em \RB{8.889}{7.970}{0.503}{%
\RG{0.000}{-0.003}{10.9091}{bb}{75}%
}\hskip.1em\allowbreak{} 为数域，且\allowbreak{}\hskip.1em \RB{8.366}{8.308}{0.503}{%
\RG{0.000}{-0.003}{10.9091}{mi}{65}%
}\hskip.1em\allowbreak{} 的特征多项式等于极小多项式。证明：矩阵方程\allowbreak{}\hskip.1em \RB{56.869}{9.932}{0.503}{%
\RG{0.000}{-0.003}{10.9091}{mi}{88}\RG{9.884}{-0.003}{10.9091}{mi}{65}\RG{21.106}{-0.003}{10.9091}{rm}{61}\RG{32.621}{-0.003}{10.9091}{mi}{65}\RG{40.803}{3.956}{7.9701}{rm}{84}\RG{47.085}{-0.003}{10.9091}{mi}{88}%
}\hskip.1em\allowbreak{}的每个解\allowbreak{}\hskip.1em \RB{55.827}{8.684}{3.225}{%
\RG{0.000}{0.002}{10.9091}{mi}{88}\RG{12.924}{0.002}{10.9091}{sy}{50}\RG{23.227}{0.002}{10.9091}{mi}{77}\RG{33.811}{-1.634}{7.9701}{mi}{110}\RG{39.447}{0.002}{10.9091}{rm}{40}\RG{43.689}{0.002}{10.9091}{bb}{75}\RG{52.174}{0.002}{10.9091}{rm}{41}%
}\hskip.1em\allowbreak{} 都是对称阵。证明：先说明题设保证存在一个循环向量\allowbreak{}\hskip.1em \RB{32.876}{7.972}{0.938}{%
\RG{0.000}{-0.002}{10.9091}{mi}{118}\RG{8.710}{-0.002}{10.9091}{sy}{50}\RG{19.013}{-0.002}{10.9091}{bb}{75}\RG{27.498}{3.957}{7.9701}{mi}{110}%
}\hskip.1em\allowbreak{} ，即

\[
\RB{78.876}{10.309}{2.601}{%
\RG{56.813}{4.509}{7.9701}{mi}{110}\RG{61.951}{4.509}{7.9701}{sy}{0}\RG{68.539}{4.509}{7.9701}{rm}{49}\RG{0.000}{0.004}{10.9091}{mi}{118}\RG{5.662}{0.004}{10.9091}{mi}{59}\RG{10.506}{0.004}{10.9091}{mi}{65}\RG{18.687}{0.004}{10.9091}{mi}{118}\RG{24.360}{0.004}{10.9091}{mi}{59}\RG{29.204}{0.004}{10.9091}{mi}{58}\RG{34.037}{0.004}{10.9091}{mi}{58}\RG{38.880}{0.004}{10.9091}{mi}{58}\RG{43.724}{0.004}{10.9091}{mi}{59}\RG{48.557}{0.004}{10.9091}{mi}{65}\RG{73.271}{0.004}{10.9091}{mi}{118}%
}
\]
构成\allowbreak{}\hskip.1em \RB{13.863}{7.974}{0.500}{%
\RG{0.000}{-0.000}{10.9091}{bb}{75}\RG{8.485}{3.959}{7.9701}{mi}{110}%
}\hskip.1em\allowbreak{} 的一组基。

记极小多项式为\allowbreak{}\hskip.1em \RB{23.708}{8.679}{3.230}{%
\RG{0.000}{-0.003}{10.9091}{mi}{109}\RG{9.578}{-0.003}{10.9091}{rm}{40}\RG{13.820}{-0.003}{10.9091}{mi}{120}\RG{20.055}{-0.003}{10.9091}{rm}{41}%
}\hskip.1em\allowbreak{} 。对\allowbreak{}\hskip.1em \RB{9.751}{5.319}{0.623}{%
\RG{0.000}{-0.003}{10.9091}{mi}{109}%
}\hskip.1em\allowbreak{} 的每个不同的首一不可约因子\allowbreak{}\hskip.1em \RB{5.846}{5.319}{2.619}{%
\RG{0.000}{-0.003}{10.9091}{mi}{112}%
}\hskip.1em\allowbreak{} ，考虑子空间

\[
\RB{99.599}{9.775}{4.305}{%
\RG{41.330}{8.839}{10.9092}{ex}{0}\RG{92.005}{8.839}{10.9092}{ex}{1}\RG{0.000}{0.003}{10.9091}{mi}{85}\RG{15.240}{0.003}{10.9091}{rm}{61}\RG{26.738}{0.003}{10.9091}{rm}{107}\RG{32.193}{0.003}{10.9091}{rm}{101}\RG{37.037}{0.003}{10.9091}{rm}{114}\RG{46.330}{0.003}{10.9091}{rm}{40}\RG{50.573}{0.003}{10.9091}{mi}{109}\RG{60.152}{0.003}{10.9091}{mi}{61}\RG{65.606}{0.003}{10.9091}{mi}{112}\RG{71.095}{0.003}{10.9091}{rm}{41}\RG{75.327}{0.003}{10.9091}{rm}{40}\RG{79.580}{0.003}{10.9091}{mi}{65}\RG{87.762}{0.003}{10.9091}{rm}{41}\RG{97.005}{0.003}{10.9091}{mi}{58}\RG{7.448}{-1.634}{7.9701}{mi}{112}%
}
\]
\noindent \RB{465.386}{9.226}{20.895}{%
\RT{0.000}{-0.001}{10.9091}{由}\RT{10.909}{-0.001}{10.9091}{于}\RG{25.247}{-0.001}{10.9091}{rm}{100}\RG{31.302}{-0.001}{10.9091}{rm}{101}\RG{36.145}{-0.001}{10.9091}{rm}{103}\RG{41.742}{-0.001}{10.9091}{rm}{40}\RG{46.004}{-0.001}{10.9091}{mi}{109}\RG{55.583}{-0.001}{10.9091}{mi}{61}\RG{61.037}{-0.001}{10.9091}{mi}{112}\RG{66.527}{-0.001}{10.9091}{rm}{41}\RG{73.799}{-0.001}{10.9091}{mi}{60}\RG{85.315}{-0.001}{10.9091}{rm}{100}\RG{91.369}{-0.001}{10.9091}{rm}{101}\RG{96.213}{-0.001}{10.9091}{rm}{103}\RG{103.648}{-0.001}{10.9091}{mi}{109}\RT{113.226}{-0.001}{10.9091}{，}\RT{123.219}{-0.001}{10.9091}{极}\RT{134.128}{-0.001}{10.9091}{小}\RT{145.037}{-0.001}{10.9091}{多}\RT{155.946}{-0.001}{10.9091}{项}\RT{166.856}{-0.001}{10.9091}{式}\RT{177.765}{-0.001}{10.9091}{的}\RT{188.674}{-0.001}{10.9091}{定}\RT{199.583}{-0.001}{10.9091}{义}\RT{210.492}{-0.001}{10.9091}{保}\RT{221.401}{-0.001}{10.9091}{证}\RG{235.735}{-0.001}{10.9091}{rm}{40}\RG{239.977}{-0.001}{10.9091}{mi}{109}\RG{249.555}{-0.001}{10.9091}{mi}{61}\RG{255.010}{-0.001}{10.9091}{mi}{112}\RG{260.499}{-0.001}{10.9091}{rm}{41}\RG{264.732}{-0.001}{10.9091}{rm}{40}\RG{268.984}{-0.001}{10.9091}{mi}{65}\RG{277.166}{-0.001}{10.9091}{rm}{41}\RG{284.438}{-0.001}{10.9091}{sy}{54}\RG{284.438}{-0.001}{10.9091}{rm}{61}\RG{295.937}{-0.001}{10.9091}{rm}{48}\RT{301.408}{-0.001}{10.9091}{，}\RT{311.401}{-0.001}{10.9091}{故}\RG{325.735}{-0.001}{10.9091}{mi}{85}\RT{341.373}{-0.001}{10.9091}{是}\RT{352.282}{-0.001}{10.9091}{真}\RT{363.191}{-0.001}{10.9091}{子}\RT{374.100}{-0.001}{10.9091}{空}\RT{385.010}{-0.001}{10.9091}{间}\RT{395.919}{-0.001}{10.9091}{。}\RT{405.912}{-0.001}{10.9091}{数}\RT{416.821}{-0.001}{10.9091}{域}\RG{431.154}{-0.001}{10.9091}{bb}{75}\RT{443.068}{-0.001}{10.9091}{是}\RT{453.977}{-0.001}{10.9091}{无}\RG{333.184}{-1.638}{7.9701}{mi}{112}\RG{247.674}{-6.994}{10.9092}{ex}{83}\RG{164.846}{-11.218}{7.9701}{mi}{110}\RT{0.000}{-15.176}{10.9091}{限}\RT{10.909}{-15.176}{10.9091}{域}\RT{21.818}{-15.176}{10.9091}{，}\RT{32.728}{-15.176}{10.9091}{有}\RT{43.637}{-15.176}{10.9091}{限}\RT{54.546}{-15.176}{10.9091}{个}\RT{65.455}{-15.176}{10.9091}{真}\RT{76.364}{-15.176}{10.9091}{子}\RT{87.274}{-15.176}{10.9091}{空}\RT{98.183}{-15.176}{10.9091}{间}\RT{109.092}{-15.176}{10.9091}{不}\RT{120.001}{-15.176}{10.9091}{能}\RT{130.910}{-15.176}{10.9091}{覆}\RT{141.819}{-15.176}{10.9091}{盖}\RG{156.361}{-15.176}{10.9091}{bb}{75}\RT{170.482}{-15.176}{10.9091}{，}\RT{181.392}{-15.176}{10.9091}{所}\RT{192.301}{-15.176}{10.9091}{以}\RT{203.210}{-15.176}{10.9091}{可}\RT{214.119}{-15.176}{10.9091}{选}\RG{228.661}{-15.176}{10.9091}{mi}{118}\RG{237.371}{-15.176}{10.9091}{sy}{50}\RG{238.566}{-15.176}{10.9091}{mi}{61}\RG{273.187}{-15.176}{10.9091}{mi}{85}\RT{285.396}{-15.176}{10.9091}{。}\RG{280.635}{-16.813}{7.9701}{mi}{112}\RG{256.765}{-18.402}{7.9701}{mi}{112}\RG{261.028}{-18.402}{7.9701}{sy}{106}\RG{263.380}{-18.402}{7.9701}{mi}{109}%
}

为补全上述事实，对每个真子空间取一个非零线性函数，使其在该子空间上为零。这些线性函数的乘积是非零多项式；而无限域上的非零多项式不可能在所有点都为零。后一句对变元个数归纳即可证明：将多项式视为最后一个变元的多项式，先使某个非零系数不为零，再避开一元多项式的有限个根。

设\allowbreak{}\hskip.1em \RB{14.039}{5.323}{2.215}{%
\RG{0.000}{0.001}{10.9091}{mi}{109}\RG{9.578}{-1.635}{7.9701}{mi}{118}%
}\hskip.1em\allowbreak{} 是使\allowbreak{}\hskip.1em \RB{57.246}{8.683}{3.226}{%
\RG{0.000}{0.001}{10.9091}{mi}{109}\RG{9.579}{-1.635}{7.9701}{mi}{118}\RG{14.486}{0.001}{10.9091}{rm}{40}\RG{18.729}{0.001}{10.9091}{mi}{65}\RG{26.910}{0.001}{10.9091}{rm}{41}\RG{31.153}{0.001}{10.9091}{mi}{118}\RG{40.044}{0.001}{10.9091}{rm}{61}\RG{51.728}{0.001}{10.9091}{rm}{48}%
}\hskip.1em\allowbreak{} 的最低次数首一多项式。带余除法给出\allowbreak{}\hskip.1em \RB{33.691}{8.683}{3.226}{%
\RG{0.000}{0.001}{10.9091}{mi}{109}\RG{9.579}{-1.635}{7.9701}{mi}{118}\RG{17.698}{0.001}{10.9091}{sy}{106}\RG{23.940}{0.001}{10.9091}{mi}{109}%
}\hskip.1em\allowbreak{} 。若\allowbreak{}\hskip.1em \RB{39.145}{8.312}{2.855}{%
\RG{0.000}{0.001}{10.9091}{mi}{109}\RG{9.578}{-1.635}{7.9701}{mi}{118}\RG{17.698}{0.001}{10.9091}{sy}{54}\RG{17.698}{0.001}{10.9091}{rm}{61}\RG{29.394}{0.001}{10.9091}{mi}{109}%
}\hskip.1em\allowbreak{} ，则\allowbreak{}\hskip.1em \RB{9.751}{5.323}{0.619}{%
\RG{0.000}{0.001}{10.9091}{mi}{109}%
}\hskip.1em\allowbreak{}的某个不可约因子\allowbreak{}\hskip.1em \RB{5.846}{5.318}{2.620}{%
\RG{0.000}{-0.004}{10.9091}{mi}{112}%
}\hskip.1em\allowbreak{} 在\allowbreak{}\hskip.1em \RB{14.039}{5.318}{2.220}{%
\RG{0.000}{-0.004}{10.9091}{mi}{109}\RG{9.578}{-1.640}{7.9701}{mi}{118}%
}\hskip.1em\allowbreak{} 中的指数较低，从而\allowbreak{}\hskip.1em \RB{44.456}{8.678}{3.231}{%
\RG{0.000}{-0.004}{10.9091}{mi}{109}\RG{9.578}{-1.640}{7.9701}{mi}{118}\RG{17.516}{-0.004}{10.9091}{sy}{106}\RG{23.577}{-0.004}{10.9091}{mi}{109}\RG{33.155}{-0.004}{10.9091}{mi}{61}\RG{38.610}{-0.004}{10.9091}{mi}{112}%
}\hskip.1em\allowbreak{} ，这会推出\allowbreak{}\hskip.1em \RB{31.106}{7.947}{3.686}{%
\RG{0.000}{-0.004}{10.9091}{mi}{118}\RG{8.710}{-0.004}{10.9091}{sy}{50}\RG{19.013}{-0.004}{10.9091}{mi}{85}\RG{26.461}{-1.640}{7.9701}{mi}{112}%
}\hskip.1em\allowbreak{} ，矛盾。因此

\[
\RB{164.239}{8.073}{2.745}{%
\RG{0.000}{0.002}{10.9091}{mi}{109}\RG{17.516}{0.002}{10.9091}{rm}{61}\RG{29.031}{0.002}{10.9091}{mi}{109}\RG{38.610}{0.002}{10.9091}{mi}{59}\RG{65.277}{0.002}{10.9091}{rm}{100}\RG{71.331}{0.002}{10.9091}{rm}{101}\RG{76.175}{0.002}{10.9091}{rm}{103}\RG{83.610}{0.002}{10.9091}{mi}{109}\RG{101.127}{0.002}{10.9091}{rm}{61}\RG{112.625}{0.002}{10.9091}{rm}{100}\RG{118.680}{0.002}{10.9091}{rm}{101}\RG{123.523}{0.002}{10.9091}{rm}{103}\RG{130.975}{0.002}{10.9091}{mi}{109}\RG{143.584}{0.002}{10.9091}{rm}{61}\RG{155.099}{0.002}{10.9091}{mi}{110}\RG{161.645}{0.002}{10.9091}{mi}{58}\RG{9.578}{-1.634}{7.9701}{mi}{118}\RG{93.189}{-1.634}{7.9701}{mi}{118}%
}
\]
于是\allowbreak{}\hskip.1em \RB{78.876}{9.757}{2.607}{%
\RG{0.000}{-0.002}{10.9091}{mi}{118}\RG{5.662}{-0.002}{10.9091}{mi}{59}\RG{10.506}{-0.002}{10.9091}{mi}{65}\RG{18.687}{-0.002}{10.9091}{mi}{118}\RG{24.360}{-0.002}{10.9091}{mi}{59}\RG{29.204}{-0.002}{10.9091}{mi}{58}\RG{34.037}{-0.002}{10.9091}{mi}{58}\RG{38.880}{-0.002}{10.9091}{mi}{58}\RG{43.724}{-0.002}{10.9091}{mi}{59}\RG{48.557}{-0.002}{10.9091}{mi}{65}\RG{56.813}{3.957}{7.9701}{mi}{110}\RG{61.952}{3.957}{7.9701}{sy}{0}\RG{68.539}{3.957}{7.9701}{rm}{49}\RG{73.271}{-0.002}{10.9091}{mi}{118}%
}\hskip.1em\allowbreak{} 线性无关，循环向量的存在得证。

现设\allowbreak{}\hskip.1em \RB{56.869}{9.939}{0.500}{%
\RG{0.000}{0.004}{10.9091}{mi}{88}\RG{9.884}{0.004}{10.9091}{mi}{65}\RG{21.106}{0.004}{10.9091}{rm}{61}\RG{32.621}{0.004}{10.9091}{mi}{65}\RG{40.803}{3.963}{7.9701}{rm}{84}\RG{47.085}{0.004}{10.9091}{mi}{88}%
}\hskip.1em\allowbreak{} 。由归纳法，

\[
\RB{128.884}{10.481}{3.227}{%
\RG{12.424}{4.505}{7.9701}{rm}{84}\RG{22.949}{4.505}{7.9701}{mi}{105}\RG{68.846}{4.505}{7.9701}{mi}{105}\RG{0.000}{-0.000}{10.9091}{rm}{40}\RG{4.242}{-0.000}{10.9091}{mi}{65}\RG{18.706}{-0.000}{10.9091}{rm}{41}\RG{26.330}{-0.000}{10.9091}{mi}{88}\RG{39.254}{-0.000}{10.9091}{rm}{61}\RG{50.769}{-0.000}{10.9091}{mi}{88}\RG{60.653}{-0.000}{10.9091}{mi}{65}\RG{94.046}{-0.000}{10.9091}{rm}{40}\RG{98.288}{-0.000}{10.9091}{mi}{105}\RG{105.077}{-0.000}{10.9091}{sy}{21}\RG{116.592}{-0.000}{10.9091}{rm}{48}\RG{122.047}{-0.000}{10.9091}{rm}{41}\RG{126.289}{-0.000}{10.9091}{mi}{58}%
}
\]
令\allowbreak{}\hskip.1em \RB{109.910}{9.755}{3.231}{%
\RG{0.000}{-0.004}{10.9091}{mi}{80}\RG{11.549}{-0.004}{10.9091}{rm}{61}\RG{23.047}{-0.004}{10.9091}{rm}{40}\RG{27.307}{-0.004}{10.9091}{mi}{118}\RG{32.969}{-0.004}{10.9091}{mi}{59}\RG{37.813}{-0.004}{10.9091}{mi}{65}\RG{45.995}{-0.004}{10.9091}{mi}{118}\RG{51.667}{-0.004}{10.9091}{mi}{59}\RG{56.511}{-0.004}{10.9091}{mi}{58}\RG{61.344}{-0.004}{10.9091}{mi}{58}\RG{66.188}{-0.004}{10.9091}{mi}{58}\RG{71.031}{-0.004}{10.9091}{mi}{59}\RG{75.864}{-0.004}{10.9091}{mi}{65}\RG{84.121}{3.955}{7.9701}{mi}{110}\RG{89.259}{3.955}{7.9701}{sy}{0}\RG{95.845}{3.955}{7.9701}{rm}{49}\RG{100.578}{-0.004}{10.9091}{mi}{118}\RG{106.257}{-0.004}{10.9091}{rm}{41}%
}\hskip.1em\allowbreak{} ，则\allowbreak{}\hskip.1em \RB{8.726}{7.947}{0.504}{%
\RG{0.000}{-0.004}{10.9091}{mi}{80}%
}\hskip.1em\allowbreak{} 可逆，且对\allowbreak{}\hskip.1em \RB{54.962}{7.761}{2.740}{%
\RG{0.000}{-0.004}{10.9091}{rm}{48}\RG{8.485}{-0.004}{10.9091}{sy}{20}\RG{20.000}{-0.004}{10.9091}{mi}{105}\RG{23.753}{-0.004}{10.9091}{mi}{59}\RG{28.586}{-0.004}{10.9091}{mi}{106}\RG{36.724}{-0.004}{10.9091}{mi}{60}\RG{48.233}{-0.004}{10.9091}{mi}{110}%
}\hskip.1em\allowbreak{} ，有

\[
\RB{227.171}{10.482}{40.088}{%
\RG{12.762}{4.507}{7.9701}{rm}{84}\RG{99.929}{4.507}{7.9701}{mi}{105}\RG{113.232}{4.507}{7.9701}{rm}{84}\RG{141.832}{4.507}{7.9701}{mi}{106}\RG{0.000}{0.002}{10.9091}{rm}{40}\RG{4.243}{0.002}{10.9091}{mi}{80}\RG{19.044}{0.002}{10.9091}{mi}{88}\RG{28.928}{0.002}{10.9091}{mi}{80}\RG{37.457}{0.002}{10.9091}{rm}{41}\RG{75.990}{0.002}{10.9091}{rm}{61}\RG{87.488}{0.002}{10.9091}{rm}{40}\RG{91.747}{0.002}{10.9091}{mi}{65}\RG{103.310}{0.002}{10.9091}{mi}{118}\RG{108.989}{0.002}{10.9091}{rm}{41}\RG{119.514}{0.002}{10.9091}{mi}{88}\RG{129.408}{0.002}{10.9091}{rm}{40}\RG{133.650}{0.002}{10.9091}{mi}{65}\RG{146.214}{0.002}{10.9091}{mi}{118}\RG{151.894}{0.002}{10.9091}{rm}{41}\RG{41.699}{-1.634}{7.9701}{mi}{105}\RG{44.582}{-1.634}{7.9701}{rm}{43}\RG{51.166}{-1.634}{7.9701}{rm}{49}\RG{55.403}{-1.634}{7.9701}{mi}{59}\RG{57.755}{-1.634}{7.9701}{mi}{106}\RG{61.639}{-1.634}{7.9701}{rm}{43}\RG{68.223}{-1.634}{7.9701}{rm}{49}\RG{93.184}{-13.657}{7.9701}{rm}{84}\RG{111.890}{-13.657}{7.9701}{rm}{84}\RG{122.415}{-13.657}{7.9701}{mi}{105}\RG{143.872}{-13.657}{7.9701}{mi}{106}\RG{75.990}{-18.161}{10.9091}{rm}{61}\RG{87.505}{-18.161}{10.9091}{mi}{118}\RG{99.466}{-18.161}{10.9091}{rm}{40}\RG{103.708}{-18.161}{10.9091}{mi}{65}\RG{118.172}{-18.161}{10.9091}{rm}{41}\RG{125.796}{-18.161}{10.9091}{mi}{88}\RG{135.680}{-18.161}{10.9091}{mi}{65}\RG{148.254}{-18.161}{10.9091}{mi}{118}\RG{93.184}{-31.821}{7.9701}{rm}{84}\RG{117.542}{-31.821}{7.9701}{mi}{105}\RG{120.425}{-31.821}{7.9701}{rm}{43}\RG{127.011}{-31.821}{7.9701}{mi}{106}\RG{164.380}{-31.821}{7.9701}{rm}{84}\RG{75.990}{-36.325}{10.9091}{rm}{61}\RG{87.505}{-36.325}{10.9091}{mi}{118}\RG{99.466}{-36.325}{10.9091}{mi}{88}\RG{109.350}{-36.325}{10.9091}{mi}{65}\RG{131.394}{-36.325}{10.9091}{mi}{118}\RG{140.103}{-36.325}{10.9091}{rm}{61}\RG{151.601}{-36.325}{10.9091}{rm}{40}\RG{155.861}{-36.325}{10.9091}{mi}{80}\RG{170.662}{-36.325}{10.9091}{mi}{88}\RG{180.546}{-36.325}{10.9091}{mi}{80}\RG{189.076}{-36.325}{10.9091}{rm}{41}\RG{224.577}{-36.325}{10.9091}{mi}{58}\RG{193.318}{-37.962}{7.9701}{mi}{106}\RG{197.202}{-37.961}{7.9701}{rm}{43}\RG{203.785}{-37.961}{7.9701}{rm}{49}\RG{208.023}{-37.962}{7.9701}{mi}{59}\RG{210.374}{-37.962}{7.9701}{mi}{105}\RG{213.258}{-37.961}{7.9701}{rm}{43}\RG{219.841}{-37.961}{7.9701}{rm}{49}%
}
\]
故\allowbreak{}\hskip.1em \RB{33.410}{9.936}{0.500}{%
\RG{0.000}{0.001}{10.9091}{mi}{80}\RG{8.519}{3.960}{7.9701}{rm}{84}\RG{14.801}{0.001}{10.9091}{mi}{88}\RG{24.685}{0.001}{10.9091}{mi}{80}%
}\hskip.1em\allowbreak{} 对称，即\allowbreak{}\hskip.1em \RB{87.452}{9.936}{0.500}{%
\RG{0.000}{0.001}{10.9091}{mi}{80}\RG{8.519}{3.960}{7.9701}{rm}{84}\RG{14.801}{0.001}{10.9091}{mi}{88}\RG{24.695}{3.960}{7.9701}{rm}{84}\RG{30.977}{0.001}{10.9091}{mi}{80}\RG{42.526}{0.001}{10.9091}{rm}{61}\RG{54.041}{0.001}{10.9091}{mi}{80}\RG{62.561}{3.960}{7.9701}{rm}{84}\RG{68.843}{0.001}{10.9091}{mi}{88}\RG{78.726}{0.001}{10.9091}{mi}{80}%
}\hskip.1em\allowbreak{} 。左右分别乘以可逆矩阵\allowbreak{}\hskip.1em \RB{33.935}{9.936}{3.226}{%
\RG{0.000}{0.001}{10.9091}{rm}{40}\RG{4.242}{0.001}{10.9091}{mi}{80}\RG{12.761}{3.960}{7.9701}{rm}{84}\RG{19.043}{0.001}{10.9091}{rm}{41}\RG{23.286}{3.960}{7.9701}{sy}{0}\RG{29.872}{3.960}{7.9701}{rm}{49}%
}\hskip.1em\allowbreak{} 、\allowbreak{}\hskip.1em \RB{19.168}{9.760}{0.500}{%
\RG{0.000}{0.001}{10.9091}{mi}{80}\RG{8.519}{3.960}{7.9701}{sy}{0}\RG{15.105}{3.960}{7.9701}{rm}{49}%
}\hskip.1em\allowbreak{} ，便得\allowbreak{}\hskip.1em \RB{40.506}{9.936}{0.500}{%
\RG{0.000}{0.001}{10.9091}{mi}{88}\RG{9.894}{3.960}{7.9701}{rm}{84}\RG{19.207}{0.001}{10.9091}{rm}{61}\RG{30.722}{0.001}{10.9091}{mi}{88}%
}\hskip.1em\allowbreak{} 。\allowbreak{}\hskip.1em \RB{8.366}{7.996}{0.500}{%
\RG{0.000}{0.001}{10.9091}{am}{3}%
}\hskip.1em\allowbreak{}

方法小结：循环基把矩阵方程转化为一个条目只依赖于指标和\allowbreak{}\hskip.1em \RB{21.667}{7.711}{2.736}{%
\RG{0.000}{0.000}{10.9091}{mi}{105}\RG{6.055}{0.000}{10.9091}{rm}{43}\RG{16.836}{0.000}{10.9091}{mi}{106}%
}\hskip.1em\allowbreak{} 的矩阵；其对称性由此直接呈现。

题目来源：谢启鸿《复旦大学高等代数习题集（2023版）》，第116页。

% Original page 7.
\section*{2022S08}
设\allowbreak{}\hskip.1em \RB{55.327}{8.679}{3.230}{%
\RG{0.000}{-0.003}{10.9091}{mi}{86}\RG{11.818}{-0.003}{10.9091}{rm}{61}\RG{23.333}{-0.003}{10.9091}{mi}{77}\RG{33.916}{-1.639}{7.9701}{mi}{110}\RG{39.553}{-0.003}{10.9091}{rm}{40}\RG{43.795}{-0.003}{10.9091}{bb}{67}\RG{51.674}{-0.003}{10.9091}{rm}{41}%
}\hskip.1em\allowbreak{} ，线性变换\allowbreak{}\hskip.1em \RB{49.661}{7.948}{2.881}{%
\RG{0.000}{-0.003}{10.9091}{mi}{39}\RG{8.348}{-0.003}{10.9091}{rm}{58}\RG{15.015}{-0.003}{10.9091}{mi}{86}\RG{26.833}{-0.003}{10.9091}{sy}{33}\RG{40.772}{-0.003}{10.9091}{mi}{86}%
}\hskip.1em\allowbreak{} 定义为

\[
\RB{172.907}{8.685}{3.224}{%
\RG{0.000}{0.003}{10.9091}{mi}{39}\RG{7.137}{0.003}{10.9091}{rm}{40}\RG{11.379}{0.003}{10.9091}{mi}{88}\RG{21.273}{0.003}{10.9091}{rm}{41}\RG{28.528}{0.003}{10.9091}{rm}{61}\RG{40.061}{0.003}{10.9091}{mi}{65}\RG{48.243}{0.003}{10.9091}{mi}{88}\RG{60.561}{0.003}{10.9091}{sy}{0}\RG{71.471}{0.003}{10.9091}{mi}{88}\RG{81.354}{0.003}{10.9091}{mi}{65}\RG{89.536}{0.003}{10.9091}{mi}{59}\RG{116.198}{0.003}{10.9091}{mi}{65}\RG{127.426}{0.003}{10.9091}{sy}{50}\RG{137.729}{0.003}{10.9091}{mi}{77}\RG{153.948}{0.003}{10.9091}{rm}{40}\RG{158.191}{0.003}{10.9091}{bb}{67}\RG{166.070}{0.003}{10.9091}{rm}{41}\RG{170.312}{0.003}{10.9091}{mi}{58}\RG{148.312}{-1.633}{7.9701}{mi}{110}%
}
\]
证明：\allowbreak{}\hskip.1em \RB{7.253}{5.322}{2.878}{%
\RG{0.000}{-0.000}{10.9091}{mi}{39}%
}\hskip.1em\allowbreak{} 可对角化的充分必要条件是\allowbreak{}\hskip.1em \RB{8.366}{8.311}{0.500}{%
\RG{0.000}{-0.000}{10.9091}{mi}{65}%
}\hskip.1em\allowbreak{} 可对角化。

证明：先证充分性。若\allowbreak{}\hskip.1em \RB{8.366}{8.307}{0.504}{%
\RG{0.000}{-0.004}{10.9091}{mi}{65}%
}\hskip.1em\allowbreak{} 可对角化，取可逆矩阵\allowbreak{}\hskip.1em \RB{8.726}{7.947}{0.504}{%
\RG{0.000}{-0.004}{10.9091}{mi}{80}%
}\hskip.1em\allowbreak{} ，使

\[
\RB{153.535}{10.306}{3.225}{%
\RG{8.519}{4.506}{7.9701}{sy}{0}\RG{15.106}{4.506}{7.9701}{rm}{49}\RG{0.000}{0.002}{10.9091}{mi}{80}\RG{19.838}{0.002}{10.9091}{mi}{65}\RG{28.020}{0.002}{10.9091}{mi}{80}\RG{39.569}{0.002}{10.9091}{rm}{61}\RG{51.084}{0.002}{10.9091}{mi}{68}\RG{63.450}{0.002}{10.9091}{rm}{61}\RG{74.948}{0.002}{10.9091}{rm}{100}\RG{81.002}{0.002}{10.9091}{rm}{105}\RG{84.024}{0.002}{10.9091}{rm}{97}\RG{89.479}{0.002}{10.9091}{rm}{103}\RG{95.086}{0.002}{10.9091}{rm}{40}\RG{99.359}{0.002}{10.9091}{mi}{21}\RG{110.455}{0.002}{10.9091}{mi}{59}\RG{115.288}{0.002}{10.9091}{mi}{58}\RG{120.131}{0.002}{10.9091}{mi}{58}\RG{124.975}{0.002}{10.9091}{mi}{58}\RG{129.808}{0.002}{10.9091}{mi}{59}\RG{134.651}{0.002}{10.9091}{mi}{21}\RG{146.698}{0.002}{10.9091}{rm}{41}\RG{150.940}{0.002}{10.9091}{mi}{58}\RG{105.723}{-1.635}{7.9701}{rm}{49}\RG{141.061}{-1.635}{7.9701}{mi}{110}%
}
\]
令\allowbreak{}\hskip.1em \RB{73.854}{9.757}{3.765}{%
\RG{0.000}{-0.002}{10.9091}{mi}{88}\RG{9.038}{-1.639}{7.9701}{mi}{105}\RG{11.915}{-1.639}{7.9701}{mi}{106}\RG{19.333}{-0.002}{10.9091}{rm}{61}\RG{30.848}{-0.002}{10.9091}{mi}{80}\RG{39.357}{-0.002}{10.9091}{mi}{69}\RG{47.420}{-1.639}{7.9701}{mi}{105}\RG{50.298}{-1.639}{7.9701}{mi}{106}\RG{54.686}{-0.002}{10.9091}{mi}{80}\RG{63.206}{3.957}{7.9701}{sy}{0}\RG{69.792}{3.957}{7.9701}{rm}{49}%
}\hskip.1em\allowbreak{} 。由于\allowbreak{}\hskip.1em \RB{64.445}{9.757}{0.622}{%
\RG{0.000}{-0.002}{10.9091}{mi}{88}\RG{12.925}{-0.002}{10.9091}{sy}{55}\RG{12.925}{-0.002}{10.9091}{sy}{33}\RG{26.864}{-0.002}{10.9091}{mi}{80}\RG{35.373}{-0.002}{10.9091}{mi}{88}\RG{45.268}{-0.002}{10.9091}{mi}{80}\RG{53.796}{3.957}{7.9701}{sy}{0}\RG{60.382}{3.957}{7.9701}{rm}{49}%
}\hskip.1em\allowbreak{} 可逆，\allowbreak{}\hskip.1em \RB{82.648}{8.680}{3.765}{%
\RG{0.000}{-0.002}{10.9091}{mi}{88}\RG{9.037}{-1.639}{7.9701}{mi}{105}\RG{11.914}{-1.639}{7.9701}{mi}{106}\RG{19.936}{-0.002}{10.9091}{rm}{40}\RG{24.169}{-0.002}{10.9091}{rm}{49}\RG{32.663}{-0.002}{10.9091}{sy}{20}\RG{44.178}{-0.002}{10.9091}{mi}{105}\RG{47.931}{-0.002}{10.9091}{mi}{59}\RG{52.764}{-0.002}{10.9091}{mi}{106}\RG{60.933}{-0.002}{10.9091}{sy}{20}\RG{72.448}{-0.002}{10.9091}{mi}{110}\RG{78.996}{-0.002}{10.9091}{rm}{41}%
}\hskip.1em\allowbreak{} 构成\allowbreak{}\hskip.1em \RB{8.889}{7.949}{0.742}{%
\RG{0.000}{-0.002}{10.9091}{mi}{86}%
}\hskip.1em\allowbreak{} 的一组基。并且

\[
\RB{221.705}{10.308}{3.758}{%
\RG{134.634}{4.508}{7.9701}{sy}{0}\RG{141.220}{4.508}{7.9701}{rm}{49}\RG{0.000}{0.004}{10.9091}{mi}{39}\RG{7.137}{0.004}{10.9091}{rm}{40}\RG{11.379}{0.004}{10.9091}{mi}{88}\RG{27.683}{0.004}{10.9091}{rm}{41}\RG{34.938}{0.004}{10.9091}{rm}{61}\RG{46.470}{0.004}{10.9091}{mi}{80}\RG{54.989}{0.004}{10.9091}{rm}{40}\RG{59.232}{0.004}{10.9091}{mi}{68}\RG{68.549}{0.004}{10.9091}{mi}{69}\RG{86.310}{0.004}{10.9091}{sy}{0}\RG{97.219}{0.004}{10.9091}{mi}{69}\RG{112.537}{0.004}{10.9091}{mi}{68}\RG{121.872}{0.004}{10.9091}{rm}{41}\RG{126.115}{0.004}{10.9091}{mi}{80}\RG{148.983}{0.004}{10.9091}{rm}{61}\RG{160.481}{0.004}{10.9091}{rm}{40}\RG{164.740}{0.004}{10.9091}{mi}{21}\RG{176.909}{0.004}{10.9091}{sy}{0}\RG{187.820}{0.004}{10.9091}{mi}{21}\RG{198.565}{0.004}{10.9091}{rm}{41}\RG{202.808}{0.004}{10.9091}{mi}{88}\RG{219.111}{0.004}{10.9091}{mi}{58}\RG{20.417}{-1.632}{7.9701}{mi}{105}\RG{23.294}{-1.632}{7.9701}{mi}{106}\RG{76.621}{-1.632}{7.9701}{mi}{105}\RG{79.498}{-1.632}{7.9701}{mi}{106}\RG{105.272}{-1.632}{7.9701}{mi}{105}\RG{108.149}{-1.632}{7.9701}{mi}{106}\RG{171.104}{-1.632}{7.9701}{mi}{105}\RG{194.183}{-1.632}{7.9701}{mi}{106}\RG{211.846}{-1.632}{7.9701}{mi}{105}\RG{214.723}{-1.632}{7.9701}{mi}{106}%
}
\]
故\allowbreak{}\hskip.1em \RB{7.253}{5.322}{2.878}{%
\RG{0.000}{-0.000}{10.9091}{mi}{39}%
}\hskip.1em\allowbreak{} 有一组特征向量基，从而可对角化。

再证必要性。先用到一个基本事实：可对角化线性变换在其不变子空间上的限制仍可对角化。为说明这一点，设\allowbreak{}\hskip.1em \RB{8.180}{7.886}{0.500}{%
\RG{0.000}{0.001}{10.9091}{mi}{84}%
}\hskip.1em\allowbreak{} 的不同特征值为\allowbreak{}\hskip.1em \RB{43.400}{8.192}{2.604}{%
\RG{0.000}{0.001}{10.9091}{mi}{18}\RG{5.121}{-1.635}{7.9701}{rm}{49}\RG{9.853}{0.001}{10.9091}{mi}{59}\RG{14.686}{0.001}{10.9091}{mi}{58}\RG{19.530}{0.001}{10.9091}{mi}{58}\RG{24.373}{0.001}{10.9091}{mi}{58}\RG{29.206}{0.001}{10.9091}{mi}{59}\RG{34.050}{0.001}{10.9091}{mi}{18}\RG{39.217}{-1.635}{7.9701}{mi}{114}%
}\hskip.1em\allowbreak{} ，定义

\[
\RB{99.684}{15.570}{15.568}{%
\RG{42.623}{10.363}{10.9092}{ex}{89}\RG{59.941}{7.379}{10.9091}{mi}{84}\RG{70.255}{7.379}{10.9091}{sy}{0}\RG{81.165}{7.379}{10.9091}{mi}{18}\RG{90.243}{7.379}{10.9091}{mi}{73}\RG{86.286}{5.602}{7.9701}{mi}{96}\RG{0.000}{-0.001}{10.9091}{mi}{25}\RG{11.339}{-0.001}{10.9091}{rm}{40}\RG{15.581}{-0.001}{10.9091}{mi}{84}\RG{23.471}{-0.001}{10.9091}{rm}{41}\RG{30.726}{-0.001}{10.9091}{rm}{61}\RG{97.090}{-0.001}{10.9091}{mi}{58}\RG{6.219}{-1.778}{7.9701}{mi}{107}\RG{61.591}{-7.484}{10.9091}{mi}{18}\RG{74.257}{-7.484}{10.9091}{sy}{0}\RG{85.166}{-7.484}{10.9091}{mi}{18}\RG{66.713}{-9.261}{7.9701}{mi}{107}\RG{90.287}{-9.261}{7.9701}{mi}{96}\RG{42.259}{-13.354}{7.9701}{mi}{96}\RG{45.718}{-13.354}{7.9701}{sy}{54}\RG{45.718}{-13.354}{7.9701}{rm}{61}\RG{52.305}{-13.354}{7.9701}{mi}{107}\RR{59.941}{2.726}{35.954}{0.436}%
}
\]
在\allowbreak{}\hskip.1em \RB{8.180}{7.888}{0.500}{%
\RG{0.000}{0.002}{10.9091}{mi}{84}%
}\hskip.1em\allowbreak{} 的特征向量基上验证可知，\allowbreak{}\hskip.1em \RB{27.124}{8.684}{3.225}{%
\RG{0.000}{0.002}{10.9091}{mi}{25}\RG{6.219}{-1.775}{7.9701}{mi}{107}\RG{11.338}{0.002}{10.9091}{rm}{40}\RG{15.581}{0.002}{10.9091}{mi}{84}\RG{23.471}{0.002}{10.9091}{rm}{41}%
}\hskip.1em\allowbreak{} 是到\allowbreak{}\hskip.1em \RB{9.869}{8.193}{2.355}{%
\RG{0.000}{0.002}{10.9091}{mi}{18}\RG{5.121}{-1.775}{7.9701}{mi}{107}%
}\hskip.1em\allowbreak{} 特征子空间的投影，且\allowbreak{}\hskip.1em \RB{67.490}{8.684}{3.804}{%
\RG{0.000}{8.184}{10.9092}{ex}{80}\RG{11.515}{-3.224}{7.9701}{mi}{107}\RG{18.453}{0.002}{10.9091}{mi}{25}\RG{24.671}{-1.775}{7.9701}{mi}{107}\RG{29.791}{0.002}{10.9091}{rm}{40}\RG{34.033}{0.002}{10.9091}{mi}{84}\RG{41.924}{0.002}{10.9091}{rm}{41}\RG{49.604}{0.002}{10.9091}{rm}{61}\RG{61.569}{0.002}{10.9091}{mi}{73}%
}\hskip.1em\allowbreak{} 。若\allowbreak{}\hskip.1em \RB{11.933}{7.953}{0.738}{%
\RG{0.000}{0.002}{10.9091}{mi}{87}%
}\hskip.1em\allowbreak{} 在\allowbreak{}\hskip.1em \RB{8.180}{7.883}{0.503}{%
\RG{0.000}{-0.003}{10.9091}{mi}{84}%
}\hskip.1em\allowbreak{} 下不变，则在每个\allowbreak{}\hskip.1em \RB{27.124}{8.679}{3.230}{%
\RG{0.000}{-0.003}{10.9091}{mi}{25}\RG{6.219}{-1.780}{7.9701}{mi}{107}\RG{11.339}{-0.003}{10.9091}{rm}{40}\RG{15.581}{-0.003}{10.9091}{mi}{84}\RG{23.471}{-0.003}{10.9091}{rm}{41}%
}\hskip.1em\allowbreak{} 下也不变，于是

\[
\RB{139.561}{17.647}{13.938}{%
\RG{32.578}{13.632}{7.9701}{mi}{114}\RG{26.364}{10.359}{10.9092}{ex}{77}\RG{42.849}{8.832}{10.9092}{ex}{0}\RG{131.966}{8.832}{10.9092}{ex}{1}\RG{0.000}{-0.004}{10.9091}{mi}{87}\RG{14.849}{-0.004}{10.9091}{rm}{61}\RG{47.849}{-0.004}{10.9091}{mi}{87}\RG{62.092}{-0.004}{10.9091}{sy}{92}\RG{71.789}{-0.004}{10.9091}{rm}{107}\RG{77.243}{-0.004}{10.9091}{rm}{101}\RG{82.087}{-0.004}{10.9091}{rm}{114}\RG{86.352}{-0.004}{10.9091}{rm}{40}\RG{90.608}{-0.004}{10.9091}{mi}{84}\RG{100.922}{-0.004}{10.9091}{sy}{0}\RG{111.831}{-0.004}{10.9091}{mi}{18}\RG{122.072}{-0.004}{10.9091}{mi}{73}\RG{127.723}{-0.004}{10.9091}{rm}{41}\RG{136.966}{-0.004}{10.9091}{mi}{58}\RG{116.952}{-1.782}{7.9701}{mi}{107}\RG{26.885}{-13.358}{7.9701}{mi}{107}\RG{31.507}{-13.358}{7.9701}{rm}{61}\RG{38.091}{-13.358}{7.9701}{rm}{49}%
}
\]
在各个直和分量中分别取基，即得\allowbreak{}\hskip.1em \RB{20.267}{8.678}{3.231}{%
\RG{0.000}{-0.004}{10.9091}{mi}{84}\RG{7.890}{-0.004}{10.9091}{sy}{106}\RG{10.920}{-1.692}{7.9701}{mi}{87}%
}\hskip.1em\allowbreak{} 的特征向量基。

现在假设\allowbreak{}\hskip.1em \RB{7.253}{5.325}{2.875}{%
\RG{0.000}{0.003}{10.9091}{mi}{39}%
}\hskip.1em\allowbreak{} 可对角化。取\allowbreak{}\hskip.1em \RB{14.158}{9.937}{0.500}{%
\RG{0.000}{0.003}{10.9091}{mi}{65}\RG{8.182}{3.961}{7.9701}{rm}{84}%
}\hskip.1em\allowbreak{} 的一个特征值\allowbreak{}\hskip.1em \RB{6.740}{5.325}{2.853}{%
\RG{0.000}{0.003}{10.9091}{mi}{22}%
}\hskip.1em\allowbreak{} 及非零特征向量\allowbreak{}\hskip.1em \RB{5.027}{8.074}{0.617}{%
\RG{0.000}{0.003}{10.9091}{mi}{98}%
}\hskip.1em\allowbreak{} ，使\allowbreak{}\hskip.1em \RB{50.922}{9.937}{2.853}{%
\RG{0.000}{0.003}{10.9091}{mi}{98}\RG{4.682}{3.961}{7.9701}{rm}{84}\RG{10.964}{0.003}{10.9091}{mi}{65}\RG{22.176}{0.003}{10.9091}{rm}{61}\RG{33.691}{0.003}{10.9091}{mi}{22}\RG{40.259}{0.003}{10.9091}{mi}{98}\RG{44.946}{3.961}{7.9701}{rm}{84}%
}\hskip.1em\allowbreak{} 。考虑子空间

\[
\RB{120.767}{10.479}{3.228}{%
\RG{42.745}{4.503}{7.9701}{rm}{84}\RG{85.576}{4.503}{7.9701}{mi}{110}\RG{0.000}{-0.001}{10.9091}{mi}{87}\RG{14.849}{-0.001}{10.9091}{rm}{61}\RG{26.364}{-0.001}{10.9091}{sy}{102}\RG{31.818}{-0.001}{10.9091}{mi}{117}\RG{38.058}{-0.001}{10.9091}{mi}{98}\RG{52.058}{-0.001}{10.9091}{rm}{58}\RG{58.118}{-0.001}{10.9091}{mi}{117}\RG{67.394}{-0.001}{10.9091}{sy}{50}\RG{77.697}{-0.001}{10.9091}{bb}{67}\RG{91.213}{-0.001}{10.9091}{sy}{103}\RG{99.689}{-0.001}{10.9091}{sy}{18}\RG{111.212}{-0.001}{10.9091}{mi}{86}\RG{118.172}{-0.001}{10.9091}{mi}{58}%
}
\]
映射\allowbreak{}\hskip.1em \RB{58.827}{8.175}{0.745}{%
\RG{0.000}{-0.005}{10.9091}{mi}{70}\RG{9.742}{-0.005}{10.9091}{rm}{58}\RG{16.409}{-0.005}{10.9091}{bb}{67}\RG{24.288}{3.954}{7.9701}{mi}{110}\RG{32.954}{-0.005}{10.9091}{sy}{33}\RG{46.894}{-0.005}{10.9091}{mi}{87}%
}\hskip.1em\allowbreak{} ，\allowbreak{}\hskip.1em \RB{54.708}{9.930}{3.232}{%
\RG{0.000}{-0.005}{10.9091}{mi}{70}\RG{8.530}{-0.005}{10.9091}{rm}{40}\RG{12.773}{-0.005}{10.9091}{mi}{117}\RG{19.018}{-0.005}{10.9091}{rm}{41}\RG{26.273}{-0.005}{10.9091}{rm}{61}\RG{37.805}{-0.005}{10.9091}{mi}{117}\RG{44.045}{-0.005}{10.9091}{mi}{98}\RG{48.732}{3.954}{7.9701}{rm}{84}%
}\hskip.1em\allowbreak{} ，为线性同构：满射由定义成立；选取\allowbreak{}\hskip.1em \RB{5.027}{8.066}{0.625}{%
\RG{0.000}{-0.005}{10.9091}{mi}{98}%
}\hskip.1em\allowbreak{} 的一个非零分量，即可看出\allowbreak{}\hskip.1em \RB{43.300}{8.682}{3.227}{%
\RG{0.000}{0.000}{10.9091}{mi}{70}\RG{8.530}{0.000}{10.9091}{rm}{40}\RG{12.773}{0.000}{10.9091}{mi}{117}\RG{19.018}{0.000}{10.9091}{rm}{41}\RG{26.273}{0.000}{10.9091}{rm}{61}\RG{37.782}{0.000}{10.9091}{rm}{48}%
}\hskip.1em\allowbreak{} 必须推出\allowbreak{}\hskip.1em \RB{26.292}{7.766}{0.740}{%
\RG{0.000}{0.000}{10.9091}{mi}{117}\RG{9.275}{0.000}{10.9091}{rm}{61}\RG{20.773}{0.000}{10.9091}{rm}{48}%
}\hskip.1em\allowbreak{} 。又有

\[
\RB{277.178}{10.476}{4.312}{%
\RG{7.137}{8.832}{10.9092}{ex}{0}\RG{35.397}{8.832}{10.9092}{ex}{1}\RG{216.114}{8.832}{10.9092}{ex}{0}\RG{269.583}{8.832}{10.9092}{ex}{1}\RG{74.052}{4.500}{7.9701}{rm}{84}\RG{104.594}{4.500}{7.9701}{rm}{84}\RG{186.756}{4.500}{7.9701}{rm}{84}\RG{0.000}{-0.004}{10.9091}{mi}{39}\RG{12.137}{-0.004}{10.9091}{mi}{70}\RG{20.667}{-0.004}{10.9091}{rm}{40}\RG{24.909}{-0.004}{10.9091}{mi}{117}\RG{31.154}{-0.004}{10.9091}{rm}{41}\RG{43.427}{-0.004}{10.9091}{rm}{61}\RG{54.942}{-0.004}{10.9091}{mi}{65}\RG{63.124}{-0.004}{10.9091}{mi}{117}\RG{69.364}{-0.004}{10.9091}{mi}{98}\RG{82.758}{-0.004}{10.9091}{sy}{0}\RG{93.667}{-0.004}{10.9091}{mi}{117}\RG{99.907}{-0.004}{10.9091}{mi}{98}\RG{110.876}{-0.004}{10.9091}{mi}{65}\RG{122.088}{-0.004}{10.9091}{rm}{61}\RG{133.586}{-0.004}{10.9091}{rm}{40}\RG{137.846}{-0.004}{10.9091}{mi}{65}\RG{148.452}{-0.004}{10.9091}{sy}{0}\RG{159.361}{-0.004}{10.9091}{mi}{22}\RG{165.929}{-0.004}{10.9091}{mi}{73}\RG{171.586}{-0.004}{10.9091}{rm}{41}\RG{175.828}{-0.004}{10.9091}{mi}{117}\RG{182.068}{-0.004}{10.9091}{mi}{98}\RG{196.068}{-0.004}{10.9091}{rm}{61}\RG{207.583}{-0.004}{10.9091}{mi}{70}\RG{221.114}{-0.004}{10.9091}{rm}{40}\RG{225.356}{-0.004}{10.9091}{mi}{65}\RG{235.962}{-0.004}{10.9091}{sy}{0}\RG{246.871}{-0.004}{10.9091}{mi}{22}\RG{253.438}{-0.004}{10.9091}{mi}{73}\RG{259.096}{-0.004}{10.9091}{rm}{41}\RG{263.338}{-0.004}{10.9091}{mi}{117}\RG{274.583}{-0.004}{10.9091}{mi}{58}%
}
\]
因此\allowbreak{}\hskip.1em \RB{11.933}{7.953}{0.738}{%
\RG{0.000}{0.002}{10.9091}{mi}{87}%
}\hskip.1em\allowbreak{} 为\allowbreak{}\hskip.1em \RB{7.253}{5.324}{2.876}{%
\RG{0.000}{0.002}{10.9091}{mi}{39}%
}\hskip.1em\allowbreak{} 的不变子空间，且\allowbreak{}\hskip.1em \RB{19.513}{8.684}{3.225}{%
\RG{0.000}{0.002}{10.9091}{mi}{39}\RG{7.136}{0.002}{10.9091}{sy}{106}\RG{10.166}{-1.687}{7.9701}{mi}{87}%
}\hskip.1em\allowbreak{} 与\allowbreak{}\hskip.1em \RB{33.850}{8.313}{2.854}{%
\RG{0.000}{0.002}{10.9091}{mi}{65}\RG{10.529}{0.002}{10.9091}{sy}{0}\RG{21.361}{0.002}{10.9091}{mi}{22}\RG{27.929}{0.002}{10.9091}{mi}{73}%
}\hskip.1em\allowbreak{} 相似。由上述基本事实，\allowbreak{}\hskip.1em \RB{33.850}{8.313}{2.854}{%
\RG{0.000}{0.002}{10.9091}{mi}{65}\RG{10.529}{0.002}{10.9091}{sy}{0}\RG{21.361}{0.002}{10.9091}{mi}{22}\RG{27.929}{0.002}{10.9091}{mi}{73}%
}\hskip.1em\allowbreak{} 可对角化；加上标量阵\allowbreak{}\hskip.1em \RB{12.489}{7.948}{2.859}{%
\RG{0.000}{-0.003}{10.9091}{mi}{22}\RG{6.567}{-0.003}{10.9091}{mi}{73}%
}\hskip.1em\allowbreak{} 不改变特征向量基，故\allowbreak{}\hskip.1em \RB{8.366}{8.308}{0.503}{%
\RG{0.000}{-0.003}{10.9091}{mi}{65}%
}\hskip.1em\allowbreak{} 也可对角化。\allowbreak{}\hskip.1em \RB{8.366}{7.992}{0.503}{%
\RG{0.000}{-0.003}{10.9091}{am}{3}%
}\hskip.1em\allowbreak{}

方法小结：固定一个左特征向量，在秩一矩阵组成的不变子空间上，交换子算子恰好等同于\allowbreak{}\hskip.1em \RB{8.366}{8.306}{0.505}{%
\RG{0.000}{-0.005}{10.9091}{mi}{65}%
}\hskip.1em\allowbreak{} 减去一个标量阵。

题目来源：谢启鸿《复旦大学高等代数习题集（2023版）》，第116页。


\end{document}
